% !TEX root = ../main.tex
Martingales are real-valued stochastic processes in continuous or discrete time that model ``fair'' betting games, in which at any time the expected future profit is always 0, irrespective of any information about the past of the game. Many proofs in probability theory are facilitated by the use of martingales. In the analysis of martingales, it will be useful to consider the process at certain random times, called \textit{stopping times} (also referred to as \textit{optional times}). We introduce these next.

\subsection{Stopping Times}
In what follows, \((\Omega, \cH, \Prob)\) is a probability space, \(\bb{T}\) a subset of \(\R\), and \(\F \idef (\F_t, t \in \bb{T})\) a filtration of sub \(\sigma\)-algebras of \(\cH\). We add a point of infinity to \(\bb{T}\) to define \(\xbar{\bb{T}} \idef \bb{T} \cup \set{\infty}\). Recall that each \(\F_t\) is a sub \(\sigma\)-algebra of \(\cH\) that models all the information on a random experiment (described by the probability space) that is available at time \(t\). If \(X\) describes our measurements on the random experiment, then the information we have available at time \(t\) is given by \(\sigma(X_s, s \leq t)\). However, \(\F_t\) could have more information than \(\sigma(X_s, s \leq t)\); that is, \(\F\) could be ``finer'' than the natural filtration of \(X\).

\begin{defn}{Stopping Time}{StoppingTime}
	A random time \(T : \Omega \to \overline{\bb{T}}\) is called a \emph{stopping time} of \(\cF\) if
	\begin{equation*}
		\{T \le t\} \in \cF_t \quad \text{for each } t \in \bb{T}.
	\end{equation*}
\end{defn}

Consider the stochastic process \((Z_t) \idef (\I_{\set{T \leq t}}, t \in \bb{T})\). Imagine that \(T\) is the time when a ``catastrophe'' happens in the random experiment, at which time the process \((Z_t)\) jumps from 0 to 1. If \(\F_t\) is the available information at time \(t\), then \(T\) is a stopping time if and only if we can tell, at time \(t\), using only the information available in \(\F_t\), whether the catastrophe has occurred yet or not. In other words, for a stopping time \(T\) we are able to construct an alarm system that uses only the information in \(\F\) and that for every \(\omega \in \Omega\) sounds exactly at the time of catastrophe \(T(\omega)\).

\begin{defn}{Past until \(T\)}{PastUntilT}
	Let \(\cF\) be a filtration on \(\overline{\bb{T}}\), and let \(T\) be a stopping time of it. The \(\sigma\)-algebra of \emph{past information until} \(T\) is defined as
	\begin{equation*}
		\cF_T \coloneqq \{H \in \cH : H \cap \{T \le t\} \in \cF_t \text{ for every } t \in \overline{\bb{T}}\}.
	\end{equation*}
\end{defn}

We can think of \(\F_T\) as the information on a random experiment that is available at a random stopping time \(T\). If \(T\) is constant, \(T = t\), then \(\F_T\) is simply \(\F_t\). As usual, we will use the notation \(\F_T\) also for the collection of numerical random variables that are \(\F_T/\xbar{\cB}\)-measurable.

\begin{thm}{\(\cF_T\) Random Variables}{FTRandomVariables}
	A random variable \(V\) belongs to \(\cF_T\) if and only if
	\begin{equation*}
		X_t \coloneqq V \I_{\{T \le t\}} \in \cF_t \text{ for every } t \in \overline{\bb{T}}.
	\end{equation*}
\end{thm}

\begin{proof}
	Without loss of generality assume \(V \ge 0\). For all \(v \ge 0\) and \(t \in \bb{T}\),
	\[
		\set{V > v} \cap \set{T \le t} = \set{X_t > v}.
	\]
	Hence, \(\set{V > v} \in \F_T\) for all \(v \ge 0\) if and only if \(X_t \in \F_t\) for every \(t \in \xbar{\bb{T}}\).

	Intuitively, \(V \in \F_T\) means that for every \(\omega\), the value \(V(\omega)\) can be discovered at time \(T(\omega)\).
\end{proof}
\subsection{Martingales}
We have already mentioned that a martingale can be thought of as a process modelling a fair game. We now make this more precise. Let \((X_t, t \in \bb{T})\) be a stochastic process describing the profit at each time \(t\) in a betting game ruled by chance. Suppose the information available at time \(t\) is given by the \(\sigma\)-algebra \(\F_t\). If for any \(s < t\) the expected incremental profit \(X_t - X_s\) is \(0\) given \(\F_s\), then the game is fair, in the sense that we cannot win whatever betting strategy we devise; the process is then a martingale.

As martingale analysis relies heavily on conditional expectations given a filtration \(\F \idef (\F_t, t \in \bb{T})\), we adopt the notation
\[
	\E_s V \;=\; \E_{\F_s} V \;=\; \E[V \mid \F_s],
\]
with \(\E_s\) being the most economical and elegant. When a filtration is not specified, the natural filtration is assumed.

\begin{defn}{(Sub/Super) Martingale}{SubSuperMartingale}
	Let \(X \coloneqq (X_t, t \in \bb{T})\) be a real-valued stochastic process that is adapted to a filtration \(\cF\), and for which each \(X_t\) is integrable. The process \(X\) is called an \(\cF\)\text{-} (sub/super) martingale if for all \(t > s\):
	\begin{equation}\label{eq:MartingaleDef}
		\begin{aligned}
			\E_s[X_t - X_s] & \ge 0 &  & \text{(submartingale)},   \\
			\E_s[X_t - X_s] & = 0   &  & \text{(martingale)},      \\
			\E_s[X_t - X_s] & \le 0 &  & \text{(supermartingale)}.
		\end{aligned}
	\end{equation}
\end{defn}

When specifying (sub/super) martingales, we will often drop the quantifier “\(\F\)-” if the corresponding filtration is obvious; for example, when it is the natural filtration of the process.

Using the betting game analogy, for submartingales the expected incremental profit is positive and for supermartingales it is negative. Think of a \textit{sub}martingale as a process that starts \textit{low} and has the tendency to \textit{increase}. You want your betting game to be a submartingale! It suffices to only consider submartingales and martingales for further analysis, as the negative of a supermartingale is a submartingale. Note that in the definition we may replace \(\E_s[X_t - X_s]\) with \(\E_s X_t - X_s\), as \(\E_s X_s = X_s\), by the “taking out what is known” property of the conditional expectation; see Theorem 4.4. From the same theorem we will frequently use the “repeated conditioning” property of the conditional expectation. In particular, taking the expectation of \(\E_s X_t - X_s\) shows that a martingale has constant expectations:

\[
	\E X_t = \E X_s.
\]

Here are some more easy properties of (sub) martingales.

\begin{thm}{Properties of (Sub) Martingales}{PropertiesOfSubMartingales}
	\begin{enumerate}
		\item When \(\bb{T} = \N\), the martingale equality \(\E_s [X_t - X_s] = 0\) holds if and only if \(\E_k [X_{k+1} - X_k] = 0, k \in \N\).
		\item If \(X\) and \(Y\) are \(\cF\)-submartingales, then so is \(aX + bY\) for \(a, b \ge 0\).
		\item If \(X\) and \(Y\) are \(\cF\)-submartingales, then so is \(X \lor Y\).
		\item Let \(f\) be a convex function on \(\R\) and \(X\) an \(\cF\)-martingale. If \(f(X_t)\) is integrable for all \(t\), then \(f(X) \coloneqq (f(X_t))\) is an \(\cF\)-submartingale.
	\end{enumerate}
\end{thm}

\subsection{Optional Stopping}
The main result of this section is Doob’s “optional” stopping theorem, which, under certain conditions, extends the martingale property, \Cref{eq:MartingaleDef} to stopping times \(S < T\), to give
\begin{equation}\label{eq:DoobsOptionalStoppingEquality}
	\E_S[X_T - X_S] = 0.
\end{equation}
and consequently (assuming \(0 \in \bb{T}\)):
\begin{equation}\label{eq:OptionalStoppingExpectation}
	\E X_T = \E X_0.
\end{equation}
This is certainly not true for \textit{all} stopping times.
\subsubsection{Stochastic Integration}

Stochastic analysis often involves the evaluation of a stochastic integral; that is, an integral of the form
\begin{equation}\label{eq:StochasticIntegral}
	\int_0^t F_s \, \di X_s,
\end{equation}
where the Integrator \(X \equiv (X_t)\) and Integrand \(F \equiv (F_t)\) are real-valued stochastic processes on a probability space \((\Omega, \cH, \Prob)\). Let \(f_t \equiv F_t(\omega)\) be the realization of \(F_t\) for a specific \(\omega \in \Omega\), and similarly let \(x_t \equiv X_t(\omega)\), with the additional assumption that \(x \equiv (x_t)\) is right-continuous and left-limited. If, for each \(\omega \in \Omega\), the real-valued function \(x\) is increasing and the function \(f \equiv (f_t)\) is measurable, then the integral can be defined pathwise; that is, for each \(\omega\) the integral
\begin{equation}\label{eq:PathwiseIntegral}
	\int_0^t f_s \, \di x_s
\end{equation}
is well-defined. Since an increasing function \(x\) corresponds to a unique measure \(\mu\), we can define this integral as the Lebesgue integral
\begin{equation}\label{eq:LebesgueMeasureIntegral}
	\int_0^t f_s \, \mu(\di s).
\end{equation}

The arguments above extend to the case when the function \(x\) can be written as the difference of two increasing functions, say \(x = y - z\). Then the integral can again be defined pathwise as
\begin{equation}\label{eq:LebesgueStieltjesDecomposition}
	\int_0^t f_s \, \di y_s \;-\; \int_0^t f_s \, \di z_s.
\end{equation}
The corresponding pathwise integral is called the Lebesgue–Stieltjes Integral.

It is not difficult to show that such functions \(x\) are of bounded variation on \([0, t]\), meaning that their total variation on \([0, t]\) is finite. The latter is defined as
\begin{equation}\label{eq:TotalVariation}
	V_x(t) \equiv \sup_{(\Pi_n)} \sum_{k=0}^{n-1} \abs{x_{s_{k+1}} - x_{s_k}},
\end{equation}
where the supremum is taken over any segmentation
\begin{equation}\label{eq:Segmentation}
	\Pi_n \equiv \set{s_0, \ldots, s_n : 0 = s_0 < s_1 < \cdots < s_n = t}, \quad n \in \N,
\end{equation}
such that the mesh of \(\Pi_n\), defined as
\begin{equation}\label{eq:MeshDefinition}
	\norm{\Pi_n} \equiv \max_{0 \le k \le n-1} (s_{k+1} - s_k),
\end{equation}
goes to \(0\) as \(n \to \infty\). Any right-continuous function of bounded variation can be written as a difference of two increasing functions.

The upshot is that for a process \(X\) that does not wiggle too much, the stochastic integral can be defined pathwise as a Lebesgue-Stieltjes integral.


\begin{defn}{Predictable Process}{PredictableProcess}
	Let \(\cF \coloneqq (\cF_t, t \ge 0)\) be a filtration. A real-valued stochastic process \(F\) is said to be \(\cF\)-\emph{predictable} (or simply \emph{predictable}) if it is measurable with respect to the \(\sigma\)-algebra
	\begin{equation}\label{eq:PredictableSigmaAlgebra}
		\cF^P \coloneqq \sigma\Bigl( \set{ H \times \left(a,b\right] : a,b \in \R_+,\, H \in \cF_a } \Bigr) \lor \sigma\Bigl( \set{ H \times \{0\} : H \in \cF_0 } \Bigr).% chktex 9
	\end{equation}
\end{defn}

Loosely speaking, this means that we are able to predict the value of \(F_t\) from all the information available before \(t\).

For discrete processes \(X \idef (X_n, n \in \N)\) and \(F \idef (F_n, n \in \N)\), the treatment of stochastic integration becomes much easier. We simply define
\begin{equation}\label{eq:DiscreteStochasticIntegral}
	\int_0^n F_k \, \di X_k \idef F_0 X_0 + F_1 (X_1 - X_0) + \cdots + F_n (X_n - X_{n-1}), \quad n = 1, 2, \ldots
\end{equation}

For \(n = 0\), the integral is defined as \(F_0 X_0\). The justification is as follows. If \(X\) is an increasing process, then for any \(\omega \in \Omega\), there corresponds to \(X(\omega)\) a discrete measure with mass \(m(k) \idef X_k(\omega) - X_{k-1}(\omega)\), at \(k = 1, \ldots, n\), and \(m(0) \idef X_0(\omega)\). The corresponding Lebesgue integral with respect to \(F(\omega)\) is then given by \Cref{eq:DiscreteStochasticIntegral}, as a function of \(\omega\). If \(X\) is arbitrary, we can write \(X\) as the difference of two increasing processes and then take the integral as the difference between the positive and negative part. This is equivalent to associating with \(X(\omega)\) a signed measure (i.e., the difference of two measures) with the same masses \(m(k)\), \(k = 0, \ldots, n\), given above. This thus defines the Lebesgue-Stieltjes integral in the discrete case; i.e., the Stieltjes integral interpreted as a Lebesgue integral.

Note that \Cref{eq:DiscreteStochasticIntegral} defines a new stochastic process
\begin{equation}\label{eq:DiscreteProcessZ}
	Z_n \idef \int_0^n F_k \, \di X_k, \quad n \in \N.
\end{equation}
Think of \(F_n\) as the number of shares owned at period \(n\) --- say the time interval from \(n-1\) to \(n\) in years --- and \(X_n\) as the price at the end of period \(n\). Then, \((X_n - X_{n-1}) F_n\) is the profit made during period \(n\). Moreover, \(Z_n\) is the total capital at the end of period \(n\), starting with an initial capital of \(X_0 F_0\). If \((X_n)\) is a martingale, then, by definition, conditional on the past until time \(n-1\), the expected increase in share price is 0; that is, \(\E_{n-1}(X_n - X_{n-1}) = 0\). We shall see that if \(F\) is a predictable process, then the expected profit during period \(n\), given the information up to period \(n-1\) is also 0. Here is the precise definition of predictability for the discrete case.


\begin{defn}{Discrete Predictable Process}{DiscretePredictableProcess}
	Given a filtration \(\cF \coloneqq (\cF_n, n \in \N)\), a process \(F \coloneqq (F_n, n \in \N)\) is said to be \emph{predictable} if \(F_0 \in \cF_0\) and \(F_n \in \cF_{n-1}\) for \(n = 1, 2, \ldots\)
\end{defn}

Continuing the financial analogy, if the shares process \(F\) is predictable with respect to the natural filtration of the share price — that is, if the number of shares owned at period \(n\) is completely determined by the evolution of the share price in the preceding periods, then \(\E_{n-1}(X_n - X_{n-1})F_n = F_n \E_{n-1}(X_n - X_{n-1}) = 0\), as reported. The total capital process \((Z_n)\) defined in \Cref{eq:DiscreteProcessZ} is in fact a martingale, as the following theorem shows:

\begin{thm}{Integral Transform of a Martingale}{IntegralTransformOfAMartingale}
	Let \(F\) be a bounded predictable process and \(X\) a martingale adapted to a filtration \(\cF\). Then, the integral transform \(Z_n \coloneqq \int_0^n F_k \di X_k, n \in \N\) is an \(\cF\)-martingale as well.
\end{thm}

It is sometimes useful to consider a stochastic process that is stopped at some random time \(T\).

\begin{defn}{Stopped Stochastic Process}{StoppedStochasticProcess}
	Let \(X \coloneqq (X_t, t \in \bb{T})\) be a stochastic process and \(T\) a random time with values in \(\overline{\bb{T}}\). By the process \(X\) \emph{stopped at time \(T\)} we mean the process \((X^\dagger_t, t \in \bb{T})\) defined by
	\begin{equation}\label{eq:StoppedProcessDefinition}
		X^\dagger_t(\omega) \coloneqq X_{T(\omega) \land t}(\omega) =
		\begin{cases}
			X_t(\omega)           & \text{if } t \le T(\omega), \\
			X_{T(\omega)}(\omega) & \text{if } t > T(\omega).
		\end{cases}
	\end{equation}
\end{defn}

Stopping a martingale at a random stopping time, gives again a martingale.

\begin{thm}{Stopped Martingale}{StoppedMartingale}
	Let \(M \coloneqq (M_n, n \in \N)\) be a martingale and \(T\) a stopping time with values in \(\overline{\N}\). Then, the process \(M\) stopped at time \(T\) is again a martingale.
\end{thm}

\begin{figure}[H]
	\hfuzz=20pt\relax
	\centering
	%% Creator: Matplotlib, PGF backend
%%
%% To include the figure in your LaTeX document, write
%%   \input{<filename>.pgf}
%%
%% Make sure the required packages are loaded in your preamble
%%   \usepackage{pgf}
%%
%% Also ensure that all the required font packages are loaded; for instance,
%% the lmodern package is sometimes necessary when using math font.
%%   \usepackage{lmodern}
%%
%% Figures using additional raster images can only be included by \input if
%% they are in the same directory as the main LaTeX file. For loading figures
%% from other directories you can use the `import` package
%%   \usepackage{import}
%%
%% and then include the figures with
%%   \import{<path to file>}{<filename>.pgf}
%%
%% Matplotlib used the following preamble
%%   \def\mathdefault#1{#1}
%%   \everymath=\expandafter{\the\everymath\displaystyle}
%%   \IfFileExists{scrextend.sty}{
%%     \usepackage[fontsize=10.000000pt]{scrextend}
%%   }{
%%     \renewcommand{\normalsize}{\fontsize{10.000000}{12.000000}\selectfont}
%%     \normalsize
%%   }
%%   \usepackage{amsmath}
%%   \usepackage{amssymb}
%%   \usepackage{amsfonts}
%%   \usepackage{libertine}
%%   \usepackage[libertine]{newtxmath}
%%   \makeatletter\@ifpackageloaded{underscore}{}{\usepackage[strings]{underscore}}\makeatother
%%
\begingroup%
\makeatletter%
\begin{pgfpicture}%
\pgfpathrectangle{\pgfpointorigin}{\pgfqpoint{6.261893in}{4.025000in}}%
\pgfusepath{use as bounding box, clip}%
\begin{pgfscope}%
\pgfsetbuttcap%
\pgfsetmiterjoin%
\definecolor{currentfill}{rgb}{1.000000,1.000000,1.000000}%
\pgfsetfillcolor{currentfill}%
\pgfsetlinewidth{0.000000pt}%
\definecolor{currentstroke}{rgb}{1.000000,1.000000,1.000000}%
\pgfsetstrokecolor{currentstroke}%
\pgfsetdash{}{0pt}%
\pgfpathmoveto{\pgfqpoint{0.000000in}{0.000000in}}%
\pgfpathlineto{\pgfqpoint{6.261893in}{0.000000in}}%
\pgfpathlineto{\pgfqpoint{6.261893in}{4.025000in}}%
\pgfpathlineto{\pgfqpoint{0.000000in}{4.025000in}}%
\pgfpathlineto{\pgfqpoint{0.000000in}{0.000000in}}%
\pgfpathclose%
\pgfusepath{fill}%
\end{pgfscope}%
\begin{pgfscope}%
\pgfsetbuttcap%
\pgfsetmiterjoin%
\definecolor{currentfill}{rgb}{1.000000,1.000000,1.000000}%
\pgfsetfillcolor{currentfill}%
\pgfsetlinewidth{0.000000pt}%
\definecolor{currentstroke}{rgb}{0.000000,0.000000,0.000000}%
\pgfsetstrokecolor{currentstroke}%
\pgfsetstrokeopacity{0.000000}%
\pgfsetdash{}{0pt}%
\pgfpathmoveto{\pgfqpoint{0.050000in}{0.050000in}}%
\pgfpathlineto{\pgfqpoint{6.170226in}{0.050000in}}%
\pgfpathlineto{\pgfqpoint{6.170226in}{3.933333in}}%
\pgfpathlineto{\pgfqpoint{0.050000in}{3.933333in}}%
\pgfpathlineto{\pgfqpoint{0.050000in}{0.050000in}}%
\pgfpathclose%
\pgfusepath{fill}%
\end{pgfscope}%
\begin{pgfscope}%
\pgfsetbuttcap%
\pgfsetroundjoin%
\definecolor{currentfill}{rgb}{0.000000,0.000000,0.000000}%
\pgfsetfillcolor{currentfill}%
\pgfsetlinewidth{0.501875pt}%
\definecolor{currentstroke}{rgb}{0.000000,0.000000,0.000000}%
\pgfsetstrokecolor{currentstroke}%
\pgfsetdash{}{0pt}%
\pgfsys@defobject{currentmarker}{\pgfqpoint{0.000000in}{-0.041667in}}{\pgfqpoint{0.000000in}{0.041667in}}{%
\pgfpathmoveto{\pgfqpoint{0.000000in}{-0.041667in}}%
\pgfpathlineto{\pgfqpoint{0.000000in}{0.041667in}}%
\pgfusepath{stroke,fill}%
}%
\begin{pgfscope}%
\pgfsys@transformshift{1.719153in}{1.775926in}%
\pgfsys@useobject{currentmarker}{}%
\end{pgfscope}%
\end{pgfscope}%
\begin{pgfscope}%
\pgfsetbuttcap%
\pgfsetroundjoin%
\definecolor{currentfill}{rgb}{0.000000,0.000000,0.000000}%
\pgfsetfillcolor{currentfill}%
\pgfsetlinewidth{0.501875pt}%
\definecolor{currentstroke}{rgb}{0.000000,0.000000,0.000000}%
\pgfsetstrokecolor{currentstroke}%
\pgfsetdash{}{0pt}%
\pgfsys@defobject{currentmarker}{\pgfqpoint{0.000000in}{-0.041667in}}{\pgfqpoint{0.000000in}{0.041667in}}{%
\pgfpathmoveto{\pgfqpoint{0.000000in}{-0.041667in}}%
\pgfpathlineto{\pgfqpoint{0.000000in}{0.041667in}}%
\pgfusepath{stroke,fill}%
}%
\begin{pgfscope}%
\pgfsys@transformshift{1.719153in}{3.933333in}%
\pgfsys@useobject{currentmarker}{}%
\end{pgfscope}%
\end{pgfscope}%
\begin{pgfscope}%
\definecolor{textcolor}{rgb}{0.000000,0.000000,0.000000}%
\pgfsetstrokecolor{textcolor}%
\pgfsetfillcolor{textcolor}%
\pgftext[x=1.719153in,y=1.685648in,,top]{\color{textcolor}{\rmfamily\fontsize{10.000000}{12.000000}\selectfont\catcode`\^=\active\def^{\ifmmode\sp\else\^{}\fi}\catcode`\%=\active\def%{\%}$1$}}%
\end{pgfscope}%
\begin{pgfscope}%
\pgfsetbuttcap%
\pgfsetroundjoin%
\definecolor{currentfill}{rgb}{0.000000,0.000000,0.000000}%
\pgfsetfillcolor{currentfill}%
\pgfsetlinewidth{0.501875pt}%
\definecolor{currentstroke}{rgb}{0.000000,0.000000,0.000000}%
\pgfsetstrokecolor{currentstroke}%
\pgfsetdash{}{0pt}%
\pgfsys@defobject{currentmarker}{\pgfqpoint{0.000000in}{-0.041667in}}{\pgfqpoint{0.000000in}{0.041667in}}{%
\pgfpathmoveto{\pgfqpoint{0.000000in}{-0.041667in}}%
\pgfpathlineto{\pgfqpoint{0.000000in}{0.041667in}}%
\pgfusepath{stroke,fill}%
}%
\begin{pgfscope}%
\pgfsys@transformshift{1.006981in}{1.775926in}%
\pgfsys@useobject{currentmarker}{}%
\end{pgfscope}%
\end{pgfscope}%
\begin{pgfscope}%
\pgfsetbuttcap%
\pgfsetroundjoin%
\definecolor{currentfill}{rgb}{0.000000,0.000000,0.000000}%
\pgfsetfillcolor{currentfill}%
\pgfsetlinewidth{0.501875pt}%
\definecolor{currentstroke}{rgb}{0.000000,0.000000,0.000000}%
\pgfsetstrokecolor{currentstroke}%
\pgfsetdash{}{0pt}%
\pgfsys@defobject{currentmarker}{\pgfqpoint{0.000000in}{-0.041667in}}{\pgfqpoint{0.000000in}{0.041667in}}{%
\pgfpathmoveto{\pgfqpoint{0.000000in}{-0.041667in}}%
\pgfpathlineto{\pgfqpoint{0.000000in}{0.041667in}}%
\pgfusepath{stroke,fill}%
}%
\begin{pgfscope}%
\pgfsys@transformshift{1.006981in}{3.933333in}%
\pgfsys@useobject{currentmarker}{}%
\end{pgfscope}%
\end{pgfscope}%
\begin{pgfscope}%
\definecolor{textcolor}{rgb}{0.000000,0.000000,0.000000}%
\pgfsetstrokecolor{textcolor}%
\pgfsetfillcolor{textcolor}%
\pgftext[x=1.006981in,y=1.685648in,,top]{\color{textcolor}{\rmfamily\fontsize{10.000000}{12.000000}\selectfont\catcode`\^=\active\def^{\ifmmode\sp\else\^{}\fi}\catcode`\%=\active\def%{\%}$\tau$}}%
\end{pgfscope}%
\begin{pgfscope}%
\pgfsetbuttcap%
\pgfsetroundjoin%
\definecolor{currentfill}{rgb}{0.000000,0.000000,0.000000}%
\pgfsetfillcolor{currentfill}%
\pgfsetlinewidth{0.501875pt}%
\definecolor{currentstroke}{rgb}{0.000000,0.000000,0.000000}%
\pgfsetstrokecolor{currentstroke}%
\pgfsetdash{}{0pt}%
\pgfsys@defobject{currentmarker}{\pgfqpoint{0.000000in}{-0.041667in}}{\pgfqpoint{0.000000in}{0.041667in}}{%
\pgfpathmoveto{\pgfqpoint{0.000000in}{-0.041667in}}%
\pgfpathlineto{\pgfqpoint{0.000000in}{0.041667in}}%
\pgfusepath{stroke,fill}%
}%
\begin{pgfscope}%
\pgfsys@transformshift{3.110113in}{1.775926in}%
\pgfsys@useobject{currentmarker}{}%
\end{pgfscope}%
\end{pgfscope}%
\begin{pgfscope}%
\pgfsetbuttcap%
\pgfsetroundjoin%
\definecolor{currentfill}{rgb}{0.000000,0.000000,0.000000}%
\pgfsetfillcolor{currentfill}%
\pgfsetlinewidth{0.501875pt}%
\definecolor{currentstroke}{rgb}{0.000000,0.000000,0.000000}%
\pgfsetstrokecolor{currentstroke}%
\pgfsetdash{}{0pt}%
\pgfsys@defobject{currentmarker}{\pgfqpoint{0.000000in}{-0.041667in}}{\pgfqpoint{0.000000in}{0.041667in}}{%
\pgfpathmoveto{\pgfqpoint{0.000000in}{-0.041667in}}%
\pgfpathlineto{\pgfqpoint{0.000000in}{0.041667in}}%
\pgfusepath{stroke,fill}%
}%
\begin{pgfscope}%
\pgfsys@transformshift{3.110113in}{3.933333in}%
\pgfsys@useobject{currentmarker}{}%
\end{pgfscope}%
\end{pgfscope}%
\begin{pgfscope}%
\definecolor{textcolor}{rgb}{0.000000,0.000000,0.000000}%
\pgfsetstrokecolor{textcolor}%
\pgfsetfillcolor{textcolor}%
\pgftext[x=3.110113in,y=1.685648in,,top]{\color{textcolor}{\rmfamily\fontsize{10.000000}{12.000000}\selectfont\catcode`\^=\active\def^{\ifmmode\sp\else\^{}\fi}\catcode`\%=\active\def%{\%}$2$}}%
\end{pgfscope}%
\begin{pgfscope}%
\pgfsetbuttcap%
\pgfsetroundjoin%
\definecolor{currentfill}{rgb}{0.000000,0.000000,0.000000}%
\pgfsetfillcolor{currentfill}%
\pgfsetlinewidth{0.501875pt}%
\definecolor{currentstroke}{rgb}{0.000000,0.000000,0.000000}%
\pgfsetstrokecolor{currentstroke}%
\pgfsetdash{}{0pt}%
\pgfsys@defobject{currentmarker}{\pgfqpoint{0.000000in}{-0.041667in}}{\pgfqpoint{0.000000in}{0.041667in}}{%
\pgfpathmoveto{\pgfqpoint{0.000000in}{-0.041667in}}%
\pgfpathlineto{\pgfqpoint{0.000000in}{0.041667in}}%
\pgfusepath{stroke,fill}%
}%
\begin{pgfscope}%
\pgfsys@transformshift{4.501074in}{1.775926in}%
\pgfsys@useobject{currentmarker}{}%
\end{pgfscope}%
\end{pgfscope}%
\begin{pgfscope}%
\pgfsetbuttcap%
\pgfsetroundjoin%
\definecolor{currentfill}{rgb}{0.000000,0.000000,0.000000}%
\pgfsetfillcolor{currentfill}%
\pgfsetlinewidth{0.501875pt}%
\definecolor{currentstroke}{rgb}{0.000000,0.000000,0.000000}%
\pgfsetstrokecolor{currentstroke}%
\pgfsetdash{}{0pt}%
\pgfsys@defobject{currentmarker}{\pgfqpoint{0.000000in}{-0.041667in}}{\pgfqpoint{0.000000in}{0.041667in}}{%
\pgfpathmoveto{\pgfqpoint{0.000000in}{-0.041667in}}%
\pgfpathlineto{\pgfqpoint{0.000000in}{0.041667in}}%
\pgfusepath{stroke,fill}%
}%
\begin{pgfscope}%
\pgfsys@transformshift{4.501074in}{3.933333in}%
\pgfsys@useobject{currentmarker}{}%
\end{pgfscope}%
\end{pgfscope}%
\begin{pgfscope}%
\definecolor{textcolor}{rgb}{0.000000,0.000000,0.000000}%
\pgfsetstrokecolor{textcolor}%
\pgfsetfillcolor{textcolor}%
\pgftext[x=4.501074in,y=1.685648in,,top]{\color{textcolor}{\rmfamily\fontsize{10.000000}{12.000000}\selectfont\catcode`\^=\active\def^{\ifmmode\sp\else\^{}\fi}\catcode`\%=\active\def%{\%}$3$}}%
\end{pgfscope}%
\begin{pgfscope}%
\pgfsetbuttcap%
\pgfsetroundjoin%
\definecolor{currentfill}{rgb}{0.000000,0.000000,0.000000}%
\pgfsetfillcolor{currentfill}%
\pgfsetlinewidth{0.501875pt}%
\definecolor{currentstroke}{rgb}{0.000000,0.000000,0.000000}%
\pgfsetstrokecolor{currentstroke}%
\pgfsetdash{}{0pt}%
\pgfsys@defobject{currentmarker}{\pgfqpoint{0.000000in}{-0.041667in}}{\pgfqpoint{0.000000in}{0.041667in}}{%
\pgfpathmoveto{\pgfqpoint{0.000000in}{-0.041667in}}%
\pgfpathlineto{\pgfqpoint{0.000000in}{0.041667in}}%
\pgfusepath{stroke,fill}%
}%
\begin{pgfscope}%
\pgfsys@transformshift{5.892034in}{1.775926in}%
\pgfsys@useobject{currentmarker}{}%
\end{pgfscope}%
\end{pgfscope}%
\begin{pgfscope}%
\pgfsetbuttcap%
\pgfsetroundjoin%
\definecolor{currentfill}{rgb}{0.000000,0.000000,0.000000}%
\pgfsetfillcolor{currentfill}%
\pgfsetlinewidth{0.501875pt}%
\definecolor{currentstroke}{rgb}{0.000000,0.000000,0.000000}%
\pgfsetstrokecolor{currentstroke}%
\pgfsetdash{}{0pt}%
\pgfsys@defobject{currentmarker}{\pgfqpoint{0.000000in}{-0.041667in}}{\pgfqpoint{0.000000in}{0.041667in}}{%
\pgfpathmoveto{\pgfqpoint{0.000000in}{-0.041667in}}%
\pgfpathlineto{\pgfqpoint{0.000000in}{0.041667in}}%
\pgfusepath{stroke,fill}%
}%
\begin{pgfscope}%
\pgfsys@transformshift{5.892034in}{3.933333in}%
\pgfsys@useobject{currentmarker}{}%
\end{pgfscope}%
\end{pgfscope}%
\begin{pgfscope}%
\definecolor{textcolor}{rgb}{0.000000,0.000000,0.000000}%
\pgfsetstrokecolor{textcolor}%
\pgfsetfillcolor{textcolor}%
\pgftext[x=5.892034in,y=1.685648in,,top]{\color{textcolor}{\rmfamily\fontsize{10.000000}{12.000000}\selectfont\catcode`\^=\active\def^{\ifmmode\sp\else\^{}\fi}\catcode`\%=\active\def%{\%}$4$}}%
\end{pgfscope}%
\begin{pgfscope}%
\pgfsetbuttcap%
\pgfsetroundjoin%
\definecolor{currentfill}{rgb}{0.000000,0.000000,0.000000}%
\pgfsetfillcolor{currentfill}%
\pgfsetlinewidth{0.501875pt}%
\definecolor{currentstroke}{rgb}{0.000000,0.000000,0.000000}%
\pgfsetstrokecolor{currentstroke}%
\pgfsetdash{}{0pt}%
\pgfsys@defobject{currentmarker}{\pgfqpoint{0.000000in}{0.000000in}}{\pgfqpoint{0.000000in}{0.020833in}}{%
\pgfpathmoveto{\pgfqpoint{0.000000in}{0.000000in}}%
\pgfpathlineto{\pgfqpoint{0.000000in}{0.020833in}}%
\pgfusepath{stroke,fill}%
}%
\begin{pgfscope}%
\pgfsys@transformshift{0.116766in}{1.775926in}%
\pgfsys@useobject{currentmarker}{}%
\end{pgfscope}%
\end{pgfscope}%
\begin{pgfscope}%
\pgfsetbuttcap%
\pgfsetroundjoin%
\definecolor{currentfill}{rgb}{0.000000,0.000000,0.000000}%
\pgfsetfillcolor{currentfill}%
\pgfsetlinewidth{0.501875pt}%
\definecolor{currentstroke}{rgb}{0.000000,0.000000,0.000000}%
\pgfsetstrokecolor{currentstroke}%
\pgfsetdash{}{0pt}%
\pgfsys@defobject{currentmarker}{\pgfqpoint{0.000000in}{-0.020833in}}{\pgfqpoint{0.000000in}{0.000000in}}{%
\pgfpathmoveto{\pgfqpoint{0.000000in}{0.000000in}}%
\pgfpathlineto{\pgfqpoint{0.000000in}{-0.020833in}}%
\pgfusepath{stroke,fill}%
}%
\begin{pgfscope}%
\pgfsys@transformshift{0.116766in}{3.933333in}%
\pgfsys@useobject{currentmarker}{}%
\end{pgfscope}%
\end{pgfscope}%
\begin{pgfscope}%
\pgfsetbuttcap%
\pgfsetroundjoin%
\definecolor{currentfill}{rgb}{0.000000,0.000000,0.000000}%
\pgfsetfillcolor{currentfill}%
\pgfsetlinewidth{0.501875pt}%
\definecolor{currentstroke}{rgb}{0.000000,0.000000,0.000000}%
\pgfsetstrokecolor{currentstroke}%
\pgfsetdash{}{0pt}%
\pgfsys@defobject{currentmarker}{\pgfqpoint{0.000000in}{0.000000in}}{\pgfqpoint{0.000000in}{0.020833in}}{%
\pgfpathmoveto{\pgfqpoint{0.000000in}{0.000000in}}%
\pgfpathlineto{\pgfqpoint{0.000000in}{0.020833in}}%
\pgfusepath{stroke,fill}%
}%
\begin{pgfscope}%
\pgfsys@transformshift{0.294809in}{1.775926in}%
\pgfsys@useobject{currentmarker}{}%
\end{pgfscope}%
\end{pgfscope}%
\begin{pgfscope}%
\pgfsetbuttcap%
\pgfsetroundjoin%
\definecolor{currentfill}{rgb}{0.000000,0.000000,0.000000}%
\pgfsetfillcolor{currentfill}%
\pgfsetlinewidth{0.501875pt}%
\definecolor{currentstroke}{rgb}{0.000000,0.000000,0.000000}%
\pgfsetstrokecolor{currentstroke}%
\pgfsetdash{}{0pt}%
\pgfsys@defobject{currentmarker}{\pgfqpoint{0.000000in}{-0.020833in}}{\pgfqpoint{0.000000in}{0.000000in}}{%
\pgfpathmoveto{\pgfqpoint{0.000000in}{0.000000in}}%
\pgfpathlineto{\pgfqpoint{0.000000in}{-0.020833in}}%
\pgfusepath{stroke,fill}%
}%
\begin{pgfscope}%
\pgfsys@transformshift{0.294809in}{3.933333in}%
\pgfsys@useobject{currentmarker}{}%
\end{pgfscope}%
\end{pgfscope}%
\begin{pgfscope}%
\pgfsetbuttcap%
\pgfsetroundjoin%
\definecolor{currentfill}{rgb}{0.000000,0.000000,0.000000}%
\pgfsetfillcolor{currentfill}%
\pgfsetlinewidth{0.501875pt}%
\definecolor{currentstroke}{rgb}{0.000000,0.000000,0.000000}%
\pgfsetstrokecolor{currentstroke}%
\pgfsetdash{}{0pt}%
\pgfsys@defobject{currentmarker}{\pgfqpoint{0.000000in}{0.000000in}}{\pgfqpoint{0.000000in}{0.020833in}}{%
\pgfpathmoveto{\pgfqpoint{0.000000in}{0.000000in}}%
\pgfpathlineto{\pgfqpoint{0.000000in}{0.020833in}}%
\pgfusepath{stroke,fill}%
}%
\begin{pgfscope}%
\pgfsys@transformshift{0.472852in}{1.775926in}%
\pgfsys@useobject{currentmarker}{}%
\end{pgfscope}%
\end{pgfscope}%
\begin{pgfscope}%
\pgfsetbuttcap%
\pgfsetroundjoin%
\definecolor{currentfill}{rgb}{0.000000,0.000000,0.000000}%
\pgfsetfillcolor{currentfill}%
\pgfsetlinewidth{0.501875pt}%
\definecolor{currentstroke}{rgb}{0.000000,0.000000,0.000000}%
\pgfsetstrokecolor{currentstroke}%
\pgfsetdash{}{0pt}%
\pgfsys@defobject{currentmarker}{\pgfqpoint{0.000000in}{-0.020833in}}{\pgfqpoint{0.000000in}{0.000000in}}{%
\pgfpathmoveto{\pgfqpoint{0.000000in}{0.000000in}}%
\pgfpathlineto{\pgfqpoint{0.000000in}{-0.020833in}}%
\pgfusepath{stroke,fill}%
}%
\begin{pgfscope}%
\pgfsys@transformshift{0.472852in}{3.933333in}%
\pgfsys@useobject{currentmarker}{}%
\end{pgfscope}%
\end{pgfscope}%
\begin{pgfscope}%
\pgfsetbuttcap%
\pgfsetroundjoin%
\definecolor{currentfill}{rgb}{0.000000,0.000000,0.000000}%
\pgfsetfillcolor{currentfill}%
\pgfsetlinewidth{0.501875pt}%
\definecolor{currentstroke}{rgb}{0.000000,0.000000,0.000000}%
\pgfsetstrokecolor{currentstroke}%
\pgfsetdash{}{0pt}%
\pgfsys@defobject{currentmarker}{\pgfqpoint{0.000000in}{0.000000in}}{\pgfqpoint{0.000000in}{0.020833in}}{%
\pgfpathmoveto{\pgfqpoint{0.000000in}{0.000000in}}%
\pgfpathlineto{\pgfqpoint{0.000000in}{0.020833in}}%
\pgfusepath{stroke,fill}%
}%
\begin{pgfscope}%
\pgfsys@transformshift{0.650895in}{1.775926in}%
\pgfsys@useobject{currentmarker}{}%
\end{pgfscope}%
\end{pgfscope}%
\begin{pgfscope}%
\pgfsetbuttcap%
\pgfsetroundjoin%
\definecolor{currentfill}{rgb}{0.000000,0.000000,0.000000}%
\pgfsetfillcolor{currentfill}%
\pgfsetlinewidth{0.501875pt}%
\definecolor{currentstroke}{rgb}{0.000000,0.000000,0.000000}%
\pgfsetstrokecolor{currentstroke}%
\pgfsetdash{}{0pt}%
\pgfsys@defobject{currentmarker}{\pgfqpoint{0.000000in}{-0.020833in}}{\pgfqpoint{0.000000in}{0.000000in}}{%
\pgfpathmoveto{\pgfqpoint{0.000000in}{0.000000in}}%
\pgfpathlineto{\pgfqpoint{0.000000in}{-0.020833in}}%
\pgfusepath{stroke,fill}%
}%
\begin{pgfscope}%
\pgfsys@transformshift{0.650895in}{3.933333in}%
\pgfsys@useobject{currentmarker}{}%
\end{pgfscope}%
\end{pgfscope}%
\begin{pgfscope}%
\pgfsetbuttcap%
\pgfsetroundjoin%
\definecolor{currentfill}{rgb}{0.000000,0.000000,0.000000}%
\pgfsetfillcolor{currentfill}%
\pgfsetlinewidth{0.501875pt}%
\definecolor{currentstroke}{rgb}{0.000000,0.000000,0.000000}%
\pgfsetstrokecolor{currentstroke}%
\pgfsetdash{}{0pt}%
\pgfsys@defobject{currentmarker}{\pgfqpoint{0.000000in}{0.000000in}}{\pgfqpoint{0.000000in}{0.020833in}}{%
\pgfpathmoveto{\pgfqpoint{0.000000in}{0.000000in}}%
\pgfpathlineto{\pgfqpoint{0.000000in}{0.020833in}}%
\pgfusepath{stroke,fill}%
}%
\begin{pgfscope}%
\pgfsys@transformshift{0.828938in}{1.775926in}%
\pgfsys@useobject{currentmarker}{}%
\end{pgfscope}%
\end{pgfscope}%
\begin{pgfscope}%
\pgfsetbuttcap%
\pgfsetroundjoin%
\definecolor{currentfill}{rgb}{0.000000,0.000000,0.000000}%
\pgfsetfillcolor{currentfill}%
\pgfsetlinewidth{0.501875pt}%
\definecolor{currentstroke}{rgb}{0.000000,0.000000,0.000000}%
\pgfsetstrokecolor{currentstroke}%
\pgfsetdash{}{0pt}%
\pgfsys@defobject{currentmarker}{\pgfqpoint{0.000000in}{-0.020833in}}{\pgfqpoint{0.000000in}{0.000000in}}{%
\pgfpathmoveto{\pgfqpoint{0.000000in}{0.000000in}}%
\pgfpathlineto{\pgfqpoint{0.000000in}{-0.020833in}}%
\pgfusepath{stroke,fill}%
}%
\begin{pgfscope}%
\pgfsys@transformshift{0.828938in}{3.933333in}%
\pgfsys@useobject{currentmarker}{}%
\end{pgfscope}%
\end{pgfscope}%
\begin{pgfscope}%
\pgfsetbuttcap%
\pgfsetroundjoin%
\definecolor{currentfill}{rgb}{0.000000,0.000000,0.000000}%
\pgfsetfillcolor{currentfill}%
\pgfsetlinewidth{0.501875pt}%
\definecolor{currentstroke}{rgb}{0.000000,0.000000,0.000000}%
\pgfsetstrokecolor{currentstroke}%
\pgfsetdash{}{0pt}%
\pgfsys@defobject{currentmarker}{\pgfqpoint{0.000000in}{0.000000in}}{\pgfqpoint{0.000000in}{0.020833in}}{%
\pgfpathmoveto{\pgfqpoint{0.000000in}{0.000000in}}%
\pgfpathlineto{\pgfqpoint{0.000000in}{0.020833in}}%
\pgfusepath{stroke,fill}%
}%
\begin{pgfscope}%
\pgfsys@transformshift{1.185024in}{1.775926in}%
\pgfsys@useobject{currentmarker}{}%
\end{pgfscope}%
\end{pgfscope}%
\begin{pgfscope}%
\pgfsetbuttcap%
\pgfsetroundjoin%
\definecolor{currentfill}{rgb}{0.000000,0.000000,0.000000}%
\pgfsetfillcolor{currentfill}%
\pgfsetlinewidth{0.501875pt}%
\definecolor{currentstroke}{rgb}{0.000000,0.000000,0.000000}%
\pgfsetstrokecolor{currentstroke}%
\pgfsetdash{}{0pt}%
\pgfsys@defobject{currentmarker}{\pgfqpoint{0.000000in}{-0.020833in}}{\pgfqpoint{0.000000in}{0.000000in}}{%
\pgfpathmoveto{\pgfqpoint{0.000000in}{0.000000in}}%
\pgfpathlineto{\pgfqpoint{0.000000in}{-0.020833in}}%
\pgfusepath{stroke,fill}%
}%
\begin{pgfscope}%
\pgfsys@transformshift{1.185024in}{3.933333in}%
\pgfsys@useobject{currentmarker}{}%
\end{pgfscope}%
\end{pgfscope}%
\begin{pgfscope}%
\pgfsetbuttcap%
\pgfsetroundjoin%
\definecolor{currentfill}{rgb}{0.000000,0.000000,0.000000}%
\pgfsetfillcolor{currentfill}%
\pgfsetlinewidth{0.501875pt}%
\definecolor{currentstroke}{rgb}{0.000000,0.000000,0.000000}%
\pgfsetstrokecolor{currentstroke}%
\pgfsetdash{}{0pt}%
\pgfsys@defobject{currentmarker}{\pgfqpoint{0.000000in}{0.000000in}}{\pgfqpoint{0.000000in}{0.020833in}}{%
\pgfpathmoveto{\pgfqpoint{0.000000in}{0.000000in}}%
\pgfpathlineto{\pgfqpoint{0.000000in}{0.020833in}}%
\pgfusepath{stroke,fill}%
}%
\begin{pgfscope}%
\pgfsys@transformshift{1.363067in}{1.775926in}%
\pgfsys@useobject{currentmarker}{}%
\end{pgfscope}%
\end{pgfscope}%
\begin{pgfscope}%
\pgfsetbuttcap%
\pgfsetroundjoin%
\definecolor{currentfill}{rgb}{0.000000,0.000000,0.000000}%
\pgfsetfillcolor{currentfill}%
\pgfsetlinewidth{0.501875pt}%
\definecolor{currentstroke}{rgb}{0.000000,0.000000,0.000000}%
\pgfsetstrokecolor{currentstroke}%
\pgfsetdash{}{0pt}%
\pgfsys@defobject{currentmarker}{\pgfqpoint{0.000000in}{-0.020833in}}{\pgfqpoint{0.000000in}{0.000000in}}{%
\pgfpathmoveto{\pgfqpoint{0.000000in}{0.000000in}}%
\pgfpathlineto{\pgfqpoint{0.000000in}{-0.020833in}}%
\pgfusepath{stroke,fill}%
}%
\begin{pgfscope}%
\pgfsys@transformshift{1.363067in}{3.933333in}%
\pgfsys@useobject{currentmarker}{}%
\end{pgfscope}%
\end{pgfscope}%
\begin{pgfscope}%
\pgfsetbuttcap%
\pgfsetroundjoin%
\definecolor{currentfill}{rgb}{0.000000,0.000000,0.000000}%
\pgfsetfillcolor{currentfill}%
\pgfsetlinewidth{0.501875pt}%
\definecolor{currentstroke}{rgb}{0.000000,0.000000,0.000000}%
\pgfsetstrokecolor{currentstroke}%
\pgfsetdash{}{0pt}%
\pgfsys@defobject{currentmarker}{\pgfqpoint{0.000000in}{0.000000in}}{\pgfqpoint{0.000000in}{0.020833in}}{%
\pgfpathmoveto{\pgfqpoint{0.000000in}{0.000000in}}%
\pgfpathlineto{\pgfqpoint{0.000000in}{0.020833in}}%
\pgfusepath{stroke,fill}%
}%
\begin{pgfscope}%
\pgfsys@transformshift{1.541110in}{1.775926in}%
\pgfsys@useobject{currentmarker}{}%
\end{pgfscope}%
\end{pgfscope}%
\begin{pgfscope}%
\pgfsetbuttcap%
\pgfsetroundjoin%
\definecolor{currentfill}{rgb}{0.000000,0.000000,0.000000}%
\pgfsetfillcolor{currentfill}%
\pgfsetlinewidth{0.501875pt}%
\definecolor{currentstroke}{rgb}{0.000000,0.000000,0.000000}%
\pgfsetstrokecolor{currentstroke}%
\pgfsetdash{}{0pt}%
\pgfsys@defobject{currentmarker}{\pgfqpoint{0.000000in}{-0.020833in}}{\pgfqpoint{0.000000in}{0.000000in}}{%
\pgfpathmoveto{\pgfqpoint{0.000000in}{0.000000in}}%
\pgfpathlineto{\pgfqpoint{0.000000in}{-0.020833in}}%
\pgfusepath{stroke,fill}%
}%
\begin{pgfscope}%
\pgfsys@transformshift{1.541110in}{3.933333in}%
\pgfsys@useobject{currentmarker}{}%
\end{pgfscope}%
\end{pgfscope}%
\begin{pgfscope}%
\pgfsetbuttcap%
\pgfsetroundjoin%
\definecolor{currentfill}{rgb}{0.000000,0.000000,0.000000}%
\pgfsetfillcolor{currentfill}%
\pgfsetlinewidth{0.501875pt}%
\definecolor{currentstroke}{rgb}{0.000000,0.000000,0.000000}%
\pgfsetstrokecolor{currentstroke}%
\pgfsetdash{}{0pt}%
\pgfsys@defobject{currentmarker}{\pgfqpoint{0.000000in}{0.000000in}}{\pgfqpoint{0.000000in}{0.020833in}}{%
\pgfpathmoveto{\pgfqpoint{0.000000in}{0.000000in}}%
\pgfpathlineto{\pgfqpoint{0.000000in}{0.020833in}}%
\pgfusepath{stroke,fill}%
}%
\begin{pgfscope}%
\pgfsys@transformshift{1.897196in}{1.775926in}%
\pgfsys@useobject{currentmarker}{}%
\end{pgfscope}%
\end{pgfscope}%
\begin{pgfscope}%
\pgfsetbuttcap%
\pgfsetroundjoin%
\definecolor{currentfill}{rgb}{0.000000,0.000000,0.000000}%
\pgfsetfillcolor{currentfill}%
\pgfsetlinewidth{0.501875pt}%
\definecolor{currentstroke}{rgb}{0.000000,0.000000,0.000000}%
\pgfsetstrokecolor{currentstroke}%
\pgfsetdash{}{0pt}%
\pgfsys@defobject{currentmarker}{\pgfqpoint{0.000000in}{-0.020833in}}{\pgfqpoint{0.000000in}{0.000000in}}{%
\pgfpathmoveto{\pgfqpoint{0.000000in}{0.000000in}}%
\pgfpathlineto{\pgfqpoint{0.000000in}{-0.020833in}}%
\pgfusepath{stroke,fill}%
}%
\begin{pgfscope}%
\pgfsys@transformshift{1.897196in}{3.933333in}%
\pgfsys@useobject{currentmarker}{}%
\end{pgfscope}%
\end{pgfscope}%
\begin{pgfscope}%
\pgfsetbuttcap%
\pgfsetroundjoin%
\definecolor{currentfill}{rgb}{0.000000,0.000000,0.000000}%
\pgfsetfillcolor{currentfill}%
\pgfsetlinewidth{0.501875pt}%
\definecolor{currentstroke}{rgb}{0.000000,0.000000,0.000000}%
\pgfsetstrokecolor{currentstroke}%
\pgfsetdash{}{0pt}%
\pgfsys@defobject{currentmarker}{\pgfqpoint{0.000000in}{0.000000in}}{\pgfqpoint{0.000000in}{0.020833in}}{%
\pgfpathmoveto{\pgfqpoint{0.000000in}{0.000000in}}%
\pgfpathlineto{\pgfqpoint{0.000000in}{0.020833in}}%
\pgfusepath{stroke,fill}%
}%
\begin{pgfscope}%
\pgfsys@transformshift{2.075238in}{1.775926in}%
\pgfsys@useobject{currentmarker}{}%
\end{pgfscope}%
\end{pgfscope}%
\begin{pgfscope}%
\pgfsetbuttcap%
\pgfsetroundjoin%
\definecolor{currentfill}{rgb}{0.000000,0.000000,0.000000}%
\pgfsetfillcolor{currentfill}%
\pgfsetlinewidth{0.501875pt}%
\definecolor{currentstroke}{rgb}{0.000000,0.000000,0.000000}%
\pgfsetstrokecolor{currentstroke}%
\pgfsetdash{}{0pt}%
\pgfsys@defobject{currentmarker}{\pgfqpoint{0.000000in}{-0.020833in}}{\pgfqpoint{0.000000in}{0.000000in}}{%
\pgfpathmoveto{\pgfqpoint{0.000000in}{0.000000in}}%
\pgfpathlineto{\pgfqpoint{0.000000in}{-0.020833in}}%
\pgfusepath{stroke,fill}%
}%
\begin{pgfscope}%
\pgfsys@transformshift{2.075238in}{3.933333in}%
\pgfsys@useobject{currentmarker}{}%
\end{pgfscope}%
\end{pgfscope}%
\begin{pgfscope}%
\pgfsetbuttcap%
\pgfsetroundjoin%
\definecolor{currentfill}{rgb}{0.000000,0.000000,0.000000}%
\pgfsetfillcolor{currentfill}%
\pgfsetlinewidth{0.501875pt}%
\definecolor{currentstroke}{rgb}{0.000000,0.000000,0.000000}%
\pgfsetstrokecolor{currentstroke}%
\pgfsetdash{}{0pt}%
\pgfsys@defobject{currentmarker}{\pgfqpoint{0.000000in}{0.000000in}}{\pgfqpoint{0.000000in}{0.020833in}}{%
\pgfpathmoveto{\pgfqpoint{0.000000in}{0.000000in}}%
\pgfpathlineto{\pgfqpoint{0.000000in}{0.020833in}}%
\pgfusepath{stroke,fill}%
}%
\begin{pgfscope}%
\pgfsys@transformshift{2.253281in}{1.775926in}%
\pgfsys@useobject{currentmarker}{}%
\end{pgfscope}%
\end{pgfscope}%
\begin{pgfscope}%
\pgfsetbuttcap%
\pgfsetroundjoin%
\definecolor{currentfill}{rgb}{0.000000,0.000000,0.000000}%
\pgfsetfillcolor{currentfill}%
\pgfsetlinewidth{0.501875pt}%
\definecolor{currentstroke}{rgb}{0.000000,0.000000,0.000000}%
\pgfsetstrokecolor{currentstroke}%
\pgfsetdash{}{0pt}%
\pgfsys@defobject{currentmarker}{\pgfqpoint{0.000000in}{-0.020833in}}{\pgfqpoint{0.000000in}{0.000000in}}{%
\pgfpathmoveto{\pgfqpoint{0.000000in}{0.000000in}}%
\pgfpathlineto{\pgfqpoint{0.000000in}{-0.020833in}}%
\pgfusepath{stroke,fill}%
}%
\begin{pgfscope}%
\pgfsys@transformshift{2.253281in}{3.933333in}%
\pgfsys@useobject{currentmarker}{}%
\end{pgfscope}%
\end{pgfscope}%
\begin{pgfscope}%
\pgfsetbuttcap%
\pgfsetroundjoin%
\definecolor{currentfill}{rgb}{0.000000,0.000000,0.000000}%
\pgfsetfillcolor{currentfill}%
\pgfsetlinewidth{0.501875pt}%
\definecolor{currentstroke}{rgb}{0.000000,0.000000,0.000000}%
\pgfsetstrokecolor{currentstroke}%
\pgfsetdash{}{0pt}%
\pgfsys@defobject{currentmarker}{\pgfqpoint{0.000000in}{0.000000in}}{\pgfqpoint{0.000000in}{0.020833in}}{%
\pgfpathmoveto{\pgfqpoint{0.000000in}{0.000000in}}%
\pgfpathlineto{\pgfqpoint{0.000000in}{0.020833in}}%
\pgfusepath{stroke,fill}%
}%
\begin{pgfscope}%
\pgfsys@transformshift{2.431324in}{1.775926in}%
\pgfsys@useobject{currentmarker}{}%
\end{pgfscope}%
\end{pgfscope}%
\begin{pgfscope}%
\pgfsetbuttcap%
\pgfsetroundjoin%
\definecolor{currentfill}{rgb}{0.000000,0.000000,0.000000}%
\pgfsetfillcolor{currentfill}%
\pgfsetlinewidth{0.501875pt}%
\definecolor{currentstroke}{rgb}{0.000000,0.000000,0.000000}%
\pgfsetstrokecolor{currentstroke}%
\pgfsetdash{}{0pt}%
\pgfsys@defobject{currentmarker}{\pgfqpoint{0.000000in}{-0.020833in}}{\pgfqpoint{0.000000in}{0.000000in}}{%
\pgfpathmoveto{\pgfqpoint{0.000000in}{0.000000in}}%
\pgfpathlineto{\pgfqpoint{0.000000in}{-0.020833in}}%
\pgfusepath{stroke,fill}%
}%
\begin{pgfscope}%
\pgfsys@transformshift{2.431324in}{3.933333in}%
\pgfsys@useobject{currentmarker}{}%
\end{pgfscope}%
\end{pgfscope}%
\begin{pgfscope}%
\pgfsetbuttcap%
\pgfsetroundjoin%
\definecolor{currentfill}{rgb}{0.000000,0.000000,0.000000}%
\pgfsetfillcolor{currentfill}%
\pgfsetlinewidth{0.501875pt}%
\definecolor{currentstroke}{rgb}{0.000000,0.000000,0.000000}%
\pgfsetstrokecolor{currentstroke}%
\pgfsetdash{}{0pt}%
\pgfsys@defobject{currentmarker}{\pgfqpoint{0.000000in}{0.000000in}}{\pgfqpoint{0.000000in}{0.020833in}}{%
\pgfpathmoveto{\pgfqpoint{0.000000in}{0.000000in}}%
\pgfpathlineto{\pgfqpoint{0.000000in}{0.020833in}}%
\pgfusepath{stroke,fill}%
}%
\begin{pgfscope}%
\pgfsys@transformshift{2.609367in}{1.775926in}%
\pgfsys@useobject{currentmarker}{}%
\end{pgfscope}%
\end{pgfscope}%
\begin{pgfscope}%
\pgfsetbuttcap%
\pgfsetroundjoin%
\definecolor{currentfill}{rgb}{0.000000,0.000000,0.000000}%
\pgfsetfillcolor{currentfill}%
\pgfsetlinewidth{0.501875pt}%
\definecolor{currentstroke}{rgb}{0.000000,0.000000,0.000000}%
\pgfsetstrokecolor{currentstroke}%
\pgfsetdash{}{0pt}%
\pgfsys@defobject{currentmarker}{\pgfqpoint{0.000000in}{-0.020833in}}{\pgfqpoint{0.000000in}{0.000000in}}{%
\pgfpathmoveto{\pgfqpoint{0.000000in}{0.000000in}}%
\pgfpathlineto{\pgfqpoint{0.000000in}{-0.020833in}}%
\pgfusepath{stroke,fill}%
}%
\begin{pgfscope}%
\pgfsys@transformshift{2.609367in}{3.933333in}%
\pgfsys@useobject{currentmarker}{}%
\end{pgfscope}%
\end{pgfscope}%
\begin{pgfscope}%
\pgfsetbuttcap%
\pgfsetroundjoin%
\definecolor{currentfill}{rgb}{0.000000,0.000000,0.000000}%
\pgfsetfillcolor{currentfill}%
\pgfsetlinewidth{0.501875pt}%
\definecolor{currentstroke}{rgb}{0.000000,0.000000,0.000000}%
\pgfsetstrokecolor{currentstroke}%
\pgfsetdash{}{0pt}%
\pgfsys@defobject{currentmarker}{\pgfqpoint{0.000000in}{0.000000in}}{\pgfqpoint{0.000000in}{0.020833in}}{%
\pgfpathmoveto{\pgfqpoint{0.000000in}{0.000000in}}%
\pgfpathlineto{\pgfqpoint{0.000000in}{0.020833in}}%
\pgfusepath{stroke,fill}%
}%
\begin{pgfscope}%
\pgfsys@transformshift{2.787410in}{1.775926in}%
\pgfsys@useobject{currentmarker}{}%
\end{pgfscope}%
\end{pgfscope}%
\begin{pgfscope}%
\pgfsetbuttcap%
\pgfsetroundjoin%
\definecolor{currentfill}{rgb}{0.000000,0.000000,0.000000}%
\pgfsetfillcolor{currentfill}%
\pgfsetlinewidth{0.501875pt}%
\definecolor{currentstroke}{rgb}{0.000000,0.000000,0.000000}%
\pgfsetstrokecolor{currentstroke}%
\pgfsetdash{}{0pt}%
\pgfsys@defobject{currentmarker}{\pgfqpoint{0.000000in}{-0.020833in}}{\pgfqpoint{0.000000in}{0.000000in}}{%
\pgfpathmoveto{\pgfqpoint{0.000000in}{0.000000in}}%
\pgfpathlineto{\pgfqpoint{0.000000in}{-0.020833in}}%
\pgfusepath{stroke,fill}%
}%
\begin{pgfscope}%
\pgfsys@transformshift{2.787410in}{3.933333in}%
\pgfsys@useobject{currentmarker}{}%
\end{pgfscope}%
\end{pgfscope}%
\begin{pgfscope}%
\pgfsetbuttcap%
\pgfsetroundjoin%
\definecolor{currentfill}{rgb}{0.000000,0.000000,0.000000}%
\pgfsetfillcolor{currentfill}%
\pgfsetlinewidth{0.501875pt}%
\definecolor{currentstroke}{rgb}{0.000000,0.000000,0.000000}%
\pgfsetstrokecolor{currentstroke}%
\pgfsetdash{}{0pt}%
\pgfsys@defobject{currentmarker}{\pgfqpoint{0.000000in}{0.000000in}}{\pgfqpoint{0.000000in}{0.020833in}}{%
\pgfpathmoveto{\pgfqpoint{0.000000in}{0.000000in}}%
\pgfpathlineto{\pgfqpoint{0.000000in}{0.020833in}}%
\pgfusepath{stroke,fill}%
}%
\begin{pgfscope}%
\pgfsys@transformshift{2.965453in}{1.775926in}%
\pgfsys@useobject{currentmarker}{}%
\end{pgfscope}%
\end{pgfscope}%
\begin{pgfscope}%
\pgfsetbuttcap%
\pgfsetroundjoin%
\definecolor{currentfill}{rgb}{0.000000,0.000000,0.000000}%
\pgfsetfillcolor{currentfill}%
\pgfsetlinewidth{0.501875pt}%
\definecolor{currentstroke}{rgb}{0.000000,0.000000,0.000000}%
\pgfsetstrokecolor{currentstroke}%
\pgfsetdash{}{0pt}%
\pgfsys@defobject{currentmarker}{\pgfqpoint{0.000000in}{-0.020833in}}{\pgfqpoint{0.000000in}{0.000000in}}{%
\pgfpathmoveto{\pgfqpoint{0.000000in}{0.000000in}}%
\pgfpathlineto{\pgfqpoint{0.000000in}{-0.020833in}}%
\pgfusepath{stroke,fill}%
}%
\begin{pgfscope}%
\pgfsys@transformshift{2.965453in}{3.933333in}%
\pgfsys@useobject{currentmarker}{}%
\end{pgfscope}%
\end{pgfscope}%
\begin{pgfscope}%
\pgfsetbuttcap%
\pgfsetroundjoin%
\definecolor{currentfill}{rgb}{0.000000,0.000000,0.000000}%
\pgfsetfillcolor{currentfill}%
\pgfsetlinewidth{0.501875pt}%
\definecolor{currentstroke}{rgb}{0.000000,0.000000,0.000000}%
\pgfsetstrokecolor{currentstroke}%
\pgfsetdash{}{0pt}%
\pgfsys@defobject{currentmarker}{\pgfqpoint{0.000000in}{0.000000in}}{\pgfqpoint{0.000000in}{0.020833in}}{%
\pgfpathmoveto{\pgfqpoint{0.000000in}{0.000000in}}%
\pgfpathlineto{\pgfqpoint{0.000000in}{0.020833in}}%
\pgfusepath{stroke,fill}%
}%
\begin{pgfscope}%
\pgfsys@transformshift{3.143496in}{1.775926in}%
\pgfsys@useobject{currentmarker}{}%
\end{pgfscope}%
\end{pgfscope}%
\begin{pgfscope}%
\pgfsetbuttcap%
\pgfsetroundjoin%
\definecolor{currentfill}{rgb}{0.000000,0.000000,0.000000}%
\pgfsetfillcolor{currentfill}%
\pgfsetlinewidth{0.501875pt}%
\definecolor{currentstroke}{rgb}{0.000000,0.000000,0.000000}%
\pgfsetstrokecolor{currentstroke}%
\pgfsetdash{}{0pt}%
\pgfsys@defobject{currentmarker}{\pgfqpoint{0.000000in}{-0.020833in}}{\pgfqpoint{0.000000in}{0.000000in}}{%
\pgfpathmoveto{\pgfqpoint{0.000000in}{0.000000in}}%
\pgfpathlineto{\pgfqpoint{0.000000in}{-0.020833in}}%
\pgfusepath{stroke,fill}%
}%
\begin{pgfscope}%
\pgfsys@transformshift{3.143496in}{3.933333in}%
\pgfsys@useobject{currentmarker}{}%
\end{pgfscope}%
\end{pgfscope}%
\begin{pgfscope}%
\pgfsetbuttcap%
\pgfsetroundjoin%
\definecolor{currentfill}{rgb}{0.000000,0.000000,0.000000}%
\pgfsetfillcolor{currentfill}%
\pgfsetlinewidth{0.501875pt}%
\definecolor{currentstroke}{rgb}{0.000000,0.000000,0.000000}%
\pgfsetstrokecolor{currentstroke}%
\pgfsetdash{}{0pt}%
\pgfsys@defobject{currentmarker}{\pgfqpoint{0.000000in}{0.000000in}}{\pgfqpoint{0.000000in}{0.020833in}}{%
\pgfpathmoveto{\pgfqpoint{0.000000in}{0.000000in}}%
\pgfpathlineto{\pgfqpoint{0.000000in}{0.020833in}}%
\pgfusepath{stroke,fill}%
}%
\begin{pgfscope}%
\pgfsys@transformshift{3.321539in}{1.775926in}%
\pgfsys@useobject{currentmarker}{}%
\end{pgfscope}%
\end{pgfscope}%
\begin{pgfscope}%
\pgfsetbuttcap%
\pgfsetroundjoin%
\definecolor{currentfill}{rgb}{0.000000,0.000000,0.000000}%
\pgfsetfillcolor{currentfill}%
\pgfsetlinewidth{0.501875pt}%
\definecolor{currentstroke}{rgb}{0.000000,0.000000,0.000000}%
\pgfsetstrokecolor{currentstroke}%
\pgfsetdash{}{0pt}%
\pgfsys@defobject{currentmarker}{\pgfqpoint{0.000000in}{-0.020833in}}{\pgfqpoint{0.000000in}{0.000000in}}{%
\pgfpathmoveto{\pgfqpoint{0.000000in}{0.000000in}}%
\pgfpathlineto{\pgfqpoint{0.000000in}{-0.020833in}}%
\pgfusepath{stroke,fill}%
}%
\begin{pgfscope}%
\pgfsys@transformshift{3.321539in}{3.933333in}%
\pgfsys@useobject{currentmarker}{}%
\end{pgfscope}%
\end{pgfscope}%
\begin{pgfscope}%
\pgfsetbuttcap%
\pgfsetroundjoin%
\definecolor{currentfill}{rgb}{0.000000,0.000000,0.000000}%
\pgfsetfillcolor{currentfill}%
\pgfsetlinewidth{0.501875pt}%
\definecolor{currentstroke}{rgb}{0.000000,0.000000,0.000000}%
\pgfsetstrokecolor{currentstroke}%
\pgfsetdash{}{0pt}%
\pgfsys@defobject{currentmarker}{\pgfqpoint{0.000000in}{0.000000in}}{\pgfqpoint{0.000000in}{0.020833in}}{%
\pgfpathmoveto{\pgfqpoint{0.000000in}{0.000000in}}%
\pgfpathlineto{\pgfqpoint{0.000000in}{0.020833in}}%
\pgfusepath{stroke,fill}%
}%
\begin{pgfscope}%
\pgfsys@transformshift{3.499582in}{1.775926in}%
\pgfsys@useobject{currentmarker}{}%
\end{pgfscope}%
\end{pgfscope}%
\begin{pgfscope}%
\pgfsetbuttcap%
\pgfsetroundjoin%
\definecolor{currentfill}{rgb}{0.000000,0.000000,0.000000}%
\pgfsetfillcolor{currentfill}%
\pgfsetlinewidth{0.501875pt}%
\definecolor{currentstroke}{rgb}{0.000000,0.000000,0.000000}%
\pgfsetstrokecolor{currentstroke}%
\pgfsetdash{}{0pt}%
\pgfsys@defobject{currentmarker}{\pgfqpoint{0.000000in}{-0.020833in}}{\pgfqpoint{0.000000in}{0.000000in}}{%
\pgfpathmoveto{\pgfqpoint{0.000000in}{0.000000in}}%
\pgfpathlineto{\pgfqpoint{0.000000in}{-0.020833in}}%
\pgfusepath{stroke,fill}%
}%
\begin{pgfscope}%
\pgfsys@transformshift{3.499582in}{3.933333in}%
\pgfsys@useobject{currentmarker}{}%
\end{pgfscope}%
\end{pgfscope}%
\begin{pgfscope}%
\pgfsetbuttcap%
\pgfsetroundjoin%
\definecolor{currentfill}{rgb}{0.000000,0.000000,0.000000}%
\pgfsetfillcolor{currentfill}%
\pgfsetlinewidth{0.501875pt}%
\definecolor{currentstroke}{rgb}{0.000000,0.000000,0.000000}%
\pgfsetstrokecolor{currentstroke}%
\pgfsetdash{}{0pt}%
\pgfsys@defobject{currentmarker}{\pgfqpoint{0.000000in}{0.000000in}}{\pgfqpoint{0.000000in}{0.020833in}}{%
\pgfpathmoveto{\pgfqpoint{0.000000in}{0.000000in}}%
\pgfpathlineto{\pgfqpoint{0.000000in}{0.020833in}}%
\pgfusepath{stroke,fill}%
}%
\begin{pgfscope}%
\pgfsys@transformshift{3.677625in}{1.775926in}%
\pgfsys@useobject{currentmarker}{}%
\end{pgfscope}%
\end{pgfscope}%
\begin{pgfscope}%
\pgfsetbuttcap%
\pgfsetroundjoin%
\definecolor{currentfill}{rgb}{0.000000,0.000000,0.000000}%
\pgfsetfillcolor{currentfill}%
\pgfsetlinewidth{0.501875pt}%
\definecolor{currentstroke}{rgb}{0.000000,0.000000,0.000000}%
\pgfsetstrokecolor{currentstroke}%
\pgfsetdash{}{0pt}%
\pgfsys@defobject{currentmarker}{\pgfqpoint{0.000000in}{-0.020833in}}{\pgfqpoint{0.000000in}{0.000000in}}{%
\pgfpathmoveto{\pgfqpoint{0.000000in}{0.000000in}}%
\pgfpathlineto{\pgfqpoint{0.000000in}{-0.020833in}}%
\pgfusepath{stroke,fill}%
}%
\begin{pgfscope}%
\pgfsys@transformshift{3.677625in}{3.933333in}%
\pgfsys@useobject{currentmarker}{}%
\end{pgfscope}%
\end{pgfscope}%
\begin{pgfscope}%
\pgfsetbuttcap%
\pgfsetroundjoin%
\definecolor{currentfill}{rgb}{0.000000,0.000000,0.000000}%
\pgfsetfillcolor{currentfill}%
\pgfsetlinewidth{0.501875pt}%
\definecolor{currentstroke}{rgb}{0.000000,0.000000,0.000000}%
\pgfsetstrokecolor{currentstroke}%
\pgfsetdash{}{0pt}%
\pgfsys@defobject{currentmarker}{\pgfqpoint{0.000000in}{0.000000in}}{\pgfqpoint{0.000000in}{0.020833in}}{%
\pgfpathmoveto{\pgfqpoint{0.000000in}{0.000000in}}%
\pgfpathlineto{\pgfqpoint{0.000000in}{0.020833in}}%
\pgfusepath{stroke,fill}%
}%
\begin{pgfscope}%
\pgfsys@transformshift{3.855668in}{1.775926in}%
\pgfsys@useobject{currentmarker}{}%
\end{pgfscope}%
\end{pgfscope}%
\begin{pgfscope}%
\pgfsetbuttcap%
\pgfsetroundjoin%
\definecolor{currentfill}{rgb}{0.000000,0.000000,0.000000}%
\pgfsetfillcolor{currentfill}%
\pgfsetlinewidth{0.501875pt}%
\definecolor{currentstroke}{rgb}{0.000000,0.000000,0.000000}%
\pgfsetstrokecolor{currentstroke}%
\pgfsetdash{}{0pt}%
\pgfsys@defobject{currentmarker}{\pgfqpoint{0.000000in}{-0.020833in}}{\pgfqpoint{0.000000in}{0.000000in}}{%
\pgfpathmoveto{\pgfqpoint{0.000000in}{0.000000in}}%
\pgfpathlineto{\pgfqpoint{0.000000in}{-0.020833in}}%
\pgfusepath{stroke,fill}%
}%
\begin{pgfscope}%
\pgfsys@transformshift{3.855668in}{3.933333in}%
\pgfsys@useobject{currentmarker}{}%
\end{pgfscope}%
\end{pgfscope}%
\begin{pgfscope}%
\pgfsetbuttcap%
\pgfsetroundjoin%
\definecolor{currentfill}{rgb}{0.000000,0.000000,0.000000}%
\pgfsetfillcolor{currentfill}%
\pgfsetlinewidth{0.501875pt}%
\definecolor{currentstroke}{rgb}{0.000000,0.000000,0.000000}%
\pgfsetstrokecolor{currentstroke}%
\pgfsetdash{}{0pt}%
\pgfsys@defobject{currentmarker}{\pgfqpoint{0.000000in}{0.000000in}}{\pgfqpoint{0.000000in}{0.020833in}}{%
\pgfpathmoveto{\pgfqpoint{0.000000in}{0.000000in}}%
\pgfpathlineto{\pgfqpoint{0.000000in}{0.020833in}}%
\pgfusepath{stroke,fill}%
}%
\begin{pgfscope}%
\pgfsys@transformshift{4.033711in}{1.775926in}%
\pgfsys@useobject{currentmarker}{}%
\end{pgfscope}%
\end{pgfscope}%
\begin{pgfscope}%
\pgfsetbuttcap%
\pgfsetroundjoin%
\definecolor{currentfill}{rgb}{0.000000,0.000000,0.000000}%
\pgfsetfillcolor{currentfill}%
\pgfsetlinewidth{0.501875pt}%
\definecolor{currentstroke}{rgb}{0.000000,0.000000,0.000000}%
\pgfsetstrokecolor{currentstroke}%
\pgfsetdash{}{0pt}%
\pgfsys@defobject{currentmarker}{\pgfqpoint{0.000000in}{-0.020833in}}{\pgfqpoint{0.000000in}{0.000000in}}{%
\pgfpathmoveto{\pgfqpoint{0.000000in}{0.000000in}}%
\pgfpathlineto{\pgfqpoint{0.000000in}{-0.020833in}}%
\pgfusepath{stroke,fill}%
}%
\begin{pgfscope}%
\pgfsys@transformshift{4.033711in}{3.933333in}%
\pgfsys@useobject{currentmarker}{}%
\end{pgfscope}%
\end{pgfscope}%
\begin{pgfscope}%
\pgfsetbuttcap%
\pgfsetroundjoin%
\definecolor{currentfill}{rgb}{0.000000,0.000000,0.000000}%
\pgfsetfillcolor{currentfill}%
\pgfsetlinewidth{0.501875pt}%
\definecolor{currentstroke}{rgb}{0.000000,0.000000,0.000000}%
\pgfsetstrokecolor{currentstroke}%
\pgfsetdash{}{0pt}%
\pgfsys@defobject{currentmarker}{\pgfqpoint{0.000000in}{0.000000in}}{\pgfqpoint{0.000000in}{0.020833in}}{%
\pgfpathmoveto{\pgfqpoint{0.000000in}{0.000000in}}%
\pgfpathlineto{\pgfqpoint{0.000000in}{0.020833in}}%
\pgfusepath{stroke,fill}%
}%
\begin{pgfscope}%
\pgfsys@transformshift{4.211754in}{1.775926in}%
\pgfsys@useobject{currentmarker}{}%
\end{pgfscope}%
\end{pgfscope}%
\begin{pgfscope}%
\pgfsetbuttcap%
\pgfsetroundjoin%
\definecolor{currentfill}{rgb}{0.000000,0.000000,0.000000}%
\pgfsetfillcolor{currentfill}%
\pgfsetlinewidth{0.501875pt}%
\definecolor{currentstroke}{rgb}{0.000000,0.000000,0.000000}%
\pgfsetstrokecolor{currentstroke}%
\pgfsetdash{}{0pt}%
\pgfsys@defobject{currentmarker}{\pgfqpoint{0.000000in}{-0.020833in}}{\pgfqpoint{0.000000in}{0.000000in}}{%
\pgfpathmoveto{\pgfqpoint{0.000000in}{0.000000in}}%
\pgfpathlineto{\pgfqpoint{0.000000in}{-0.020833in}}%
\pgfusepath{stroke,fill}%
}%
\begin{pgfscope}%
\pgfsys@transformshift{4.211754in}{3.933333in}%
\pgfsys@useobject{currentmarker}{}%
\end{pgfscope}%
\end{pgfscope}%
\begin{pgfscope}%
\pgfsetbuttcap%
\pgfsetroundjoin%
\definecolor{currentfill}{rgb}{0.000000,0.000000,0.000000}%
\pgfsetfillcolor{currentfill}%
\pgfsetlinewidth{0.501875pt}%
\definecolor{currentstroke}{rgb}{0.000000,0.000000,0.000000}%
\pgfsetstrokecolor{currentstroke}%
\pgfsetdash{}{0pt}%
\pgfsys@defobject{currentmarker}{\pgfqpoint{0.000000in}{0.000000in}}{\pgfqpoint{0.000000in}{0.020833in}}{%
\pgfpathmoveto{\pgfqpoint{0.000000in}{0.000000in}}%
\pgfpathlineto{\pgfqpoint{0.000000in}{0.020833in}}%
\pgfusepath{stroke,fill}%
}%
\begin{pgfscope}%
\pgfsys@transformshift{4.389797in}{1.775926in}%
\pgfsys@useobject{currentmarker}{}%
\end{pgfscope}%
\end{pgfscope}%
\begin{pgfscope}%
\pgfsetbuttcap%
\pgfsetroundjoin%
\definecolor{currentfill}{rgb}{0.000000,0.000000,0.000000}%
\pgfsetfillcolor{currentfill}%
\pgfsetlinewidth{0.501875pt}%
\definecolor{currentstroke}{rgb}{0.000000,0.000000,0.000000}%
\pgfsetstrokecolor{currentstroke}%
\pgfsetdash{}{0pt}%
\pgfsys@defobject{currentmarker}{\pgfqpoint{0.000000in}{-0.020833in}}{\pgfqpoint{0.000000in}{0.000000in}}{%
\pgfpathmoveto{\pgfqpoint{0.000000in}{0.000000in}}%
\pgfpathlineto{\pgfqpoint{0.000000in}{-0.020833in}}%
\pgfusepath{stroke,fill}%
}%
\begin{pgfscope}%
\pgfsys@transformshift{4.389797in}{3.933333in}%
\pgfsys@useobject{currentmarker}{}%
\end{pgfscope}%
\end{pgfscope}%
\begin{pgfscope}%
\pgfsetbuttcap%
\pgfsetroundjoin%
\definecolor{currentfill}{rgb}{0.000000,0.000000,0.000000}%
\pgfsetfillcolor{currentfill}%
\pgfsetlinewidth{0.501875pt}%
\definecolor{currentstroke}{rgb}{0.000000,0.000000,0.000000}%
\pgfsetstrokecolor{currentstroke}%
\pgfsetdash{}{0pt}%
\pgfsys@defobject{currentmarker}{\pgfqpoint{0.000000in}{0.000000in}}{\pgfqpoint{0.000000in}{0.020833in}}{%
\pgfpathmoveto{\pgfqpoint{0.000000in}{0.000000in}}%
\pgfpathlineto{\pgfqpoint{0.000000in}{0.020833in}}%
\pgfusepath{stroke,fill}%
}%
\begin{pgfscope}%
\pgfsys@transformshift{4.567840in}{1.775926in}%
\pgfsys@useobject{currentmarker}{}%
\end{pgfscope}%
\end{pgfscope}%
\begin{pgfscope}%
\pgfsetbuttcap%
\pgfsetroundjoin%
\definecolor{currentfill}{rgb}{0.000000,0.000000,0.000000}%
\pgfsetfillcolor{currentfill}%
\pgfsetlinewidth{0.501875pt}%
\definecolor{currentstroke}{rgb}{0.000000,0.000000,0.000000}%
\pgfsetstrokecolor{currentstroke}%
\pgfsetdash{}{0pt}%
\pgfsys@defobject{currentmarker}{\pgfqpoint{0.000000in}{-0.020833in}}{\pgfqpoint{0.000000in}{0.000000in}}{%
\pgfpathmoveto{\pgfqpoint{0.000000in}{0.000000in}}%
\pgfpathlineto{\pgfqpoint{0.000000in}{-0.020833in}}%
\pgfusepath{stroke,fill}%
}%
\begin{pgfscope}%
\pgfsys@transformshift{4.567840in}{3.933333in}%
\pgfsys@useobject{currentmarker}{}%
\end{pgfscope}%
\end{pgfscope}%
\begin{pgfscope}%
\pgfsetbuttcap%
\pgfsetroundjoin%
\definecolor{currentfill}{rgb}{0.000000,0.000000,0.000000}%
\pgfsetfillcolor{currentfill}%
\pgfsetlinewidth{0.501875pt}%
\definecolor{currentstroke}{rgb}{0.000000,0.000000,0.000000}%
\pgfsetstrokecolor{currentstroke}%
\pgfsetdash{}{0pt}%
\pgfsys@defobject{currentmarker}{\pgfqpoint{0.000000in}{0.000000in}}{\pgfqpoint{0.000000in}{0.020833in}}{%
\pgfpathmoveto{\pgfqpoint{0.000000in}{0.000000in}}%
\pgfpathlineto{\pgfqpoint{0.000000in}{0.020833in}}%
\pgfusepath{stroke,fill}%
}%
\begin{pgfscope}%
\pgfsys@transformshift{4.745883in}{1.775926in}%
\pgfsys@useobject{currentmarker}{}%
\end{pgfscope}%
\end{pgfscope}%
\begin{pgfscope}%
\pgfsetbuttcap%
\pgfsetroundjoin%
\definecolor{currentfill}{rgb}{0.000000,0.000000,0.000000}%
\pgfsetfillcolor{currentfill}%
\pgfsetlinewidth{0.501875pt}%
\definecolor{currentstroke}{rgb}{0.000000,0.000000,0.000000}%
\pgfsetstrokecolor{currentstroke}%
\pgfsetdash{}{0pt}%
\pgfsys@defobject{currentmarker}{\pgfqpoint{0.000000in}{-0.020833in}}{\pgfqpoint{0.000000in}{0.000000in}}{%
\pgfpathmoveto{\pgfqpoint{0.000000in}{0.000000in}}%
\pgfpathlineto{\pgfqpoint{0.000000in}{-0.020833in}}%
\pgfusepath{stroke,fill}%
}%
\begin{pgfscope}%
\pgfsys@transformshift{4.745883in}{3.933333in}%
\pgfsys@useobject{currentmarker}{}%
\end{pgfscope}%
\end{pgfscope}%
\begin{pgfscope}%
\pgfsetbuttcap%
\pgfsetroundjoin%
\definecolor{currentfill}{rgb}{0.000000,0.000000,0.000000}%
\pgfsetfillcolor{currentfill}%
\pgfsetlinewidth{0.501875pt}%
\definecolor{currentstroke}{rgb}{0.000000,0.000000,0.000000}%
\pgfsetstrokecolor{currentstroke}%
\pgfsetdash{}{0pt}%
\pgfsys@defobject{currentmarker}{\pgfqpoint{0.000000in}{0.000000in}}{\pgfqpoint{0.000000in}{0.020833in}}{%
\pgfpathmoveto{\pgfqpoint{0.000000in}{0.000000in}}%
\pgfpathlineto{\pgfqpoint{0.000000in}{0.020833in}}%
\pgfusepath{stroke,fill}%
}%
\begin{pgfscope}%
\pgfsys@transformshift{4.923926in}{1.775926in}%
\pgfsys@useobject{currentmarker}{}%
\end{pgfscope}%
\end{pgfscope}%
\begin{pgfscope}%
\pgfsetbuttcap%
\pgfsetroundjoin%
\definecolor{currentfill}{rgb}{0.000000,0.000000,0.000000}%
\pgfsetfillcolor{currentfill}%
\pgfsetlinewidth{0.501875pt}%
\definecolor{currentstroke}{rgb}{0.000000,0.000000,0.000000}%
\pgfsetstrokecolor{currentstroke}%
\pgfsetdash{}{0pt}%
\pgfsys@defobject{currentmarker}{\pgfqpoint{0.000000in}{-0.020833in}}{\pgfqpoint{0.000000in}{0.000000in}}{%
\pgfpathmoveto{\pgfqpoint{0.000000in}{0.000000in}}%
\pgfpathlineto{\pgfqpoint{0.000000in}{-0.020833in}}%
\pgfusepath{stroke,fill}%
}%
\begin{pgfscope}%
\pgfsys@transformshift{4.923926in}{3.933333in}%
\pgfsys@useobject{currentmarker}{}%
\end{pgfscope}%
\end{pgfscope}%
\begin{pgfscope}%
\pgfsetbuttcap%
\pgfsetroundjoin%
\definecolor{currentfill}{rgb}{0.000000,0.000000,0.000000}%
\pgfsetfillcolor{currentfill}%
\pgfsetlinewidth{0.501875pt}%
\definecolor{currentstroke}{rgb}{0.000000,0.000000,0.000000}%
\pgfsetstrokecolor{currentstroke}%
\pgfsetdash{}{0pt}%
\pgfsys@defobject{currentmarker}{\pgfqpoint{0.000000in}{0.000000in}}{\pgfqpoint{0.000000in}{0.020833in}}{%
\pgfpathmoveto{\pgfqpoint{0.000000in}{0.000000in}}%
\pgfpathlineto{\pgfqpoint{0.000000in}{0.020833in}}%
\pgfusepath{stroke,fill}%
}%
\begin{pgfscope}%
\pgfsys@transformshift{5.101969in}{1.775926in}%
\pgfsys@useobject{currentmarker}{}%
\end{pgfscope}%
\end{pgfscope}%
\begin{pgfscope}%
\pgfsetbuttcap%
\pgfsetroundjoin%
\definecolor{currentfill}{rgb}{0.000000,0.000000,0.000000}%
\pgfsetfillcolor{currentfill}%
\pgfsetlinewidth{0.501875pt}%
\definecolor{currentstroke}{rgb}{0.000000,0.000000,0.000000}%
\pgfsetstrokecolor{currentstroke}%
\pgfsetdash{}{0pt}%
\pgfsys@defobject{currentmarker}{\pgfqpoint{0.000000in}{-0.020833in}}{\pgfqpoint{0.000000in}{0.000000in}}{%
\pgfpathmoveto{\pgfqpoint{0.000000in}{0.000000in}}%
\pgfpathlineto{\pgfqpoint{0.000000in}{-0.020833in}}%
\pgfusepath{stroke,fill}%
}%
\begin{pgfscope}%
\pgfsys@transformshift{5.101969in}{3.933333in}%
\pgfsys@useobject{currentmarker}{}%
\end{pgfscope}%
\end{pgfscope}%
\begin{pgfscope}%
\pgfsetbuttcap%
\pgfsetroundjoin%
\definecolor{currentfill}{rgb}{0.000000,0.000000,0.000000}%
\pgfsetfillcolor{currentfill}%
\pgfsetlinewidth{0.501875pt}%
\definecolor{currentstroke}{rgb}{0.000000,0.000000,0.000000}%
\pgfsetstrokecolor{currentstroke}%
\pgfsetdash{}{0pt}%
\pgfsys@defobject{currentmarker}{\pgfqpoint{0.000000in}{0.000000in}}{\pgfqpoint{0.000000in}{0.020833in}}{%
\pgfpathmoveto{\pgfqpoint{0.000000in}{0.000000in}}%
\pgfpathlineto{\pgfqpoint{0.000000in}{0.020833in}}%
\pgfusepath{stroke,fill}%
}%
\begin{pgfscope}%
\pgfsys@transformshift{5.280011in}{1.775926in}%
\pgfsys@useobject{currentmarker}{}%
\end{pgfscope}%
\end{pgfscope}%
\begin{pgfscope}%
\pgfsetbuttcap%
\pgfsetroundjoin%
\definecolor{currentfill}{rgb}{0.000000,0.000000,0.000000}%
\pgfsetfillcolor{currentfill}%
\pgfsetlinewidth{0.501875pt}%
\definecolor{currentstroke}{rgb}{0.000000,0.000000,0.000000}%
\pgfsetstrokecolor{currentstroke}%
\pgfsetdash{}{0pt}%
\pgfsys@defobject{currentmarker}{\pgfqpoint{0.000000in}{-0.020833in}}{\pgfqpoint{0.000000in}{0.000000in}}{%
\pgfpathmoveto{\pgfqpoint{0.000000in}{0.000000in}}%
\pgfpathlineto{\pgfqpoint{0.000000in}{-0.020833in}}%
\pgfusepath{stroke,fill}%
}%
\begin{pgfscope}%
\pgfsys@transformshift{5.280011in}{3.933333in}%
\pgfsys@useobject{currentmarker}{}%
\end{pgfscope}%
\end{pgfscope}%
\begin{pgfscope}%
\pgfsetbuttcap%
\pgfsetroundjoin%
\definecolor{currentfill}{rgb}{0.000000,0.000000,0.000000}%
\pgfsetfillcolor{currentfill}%
\pgfsetlinewidth{0.501875pt}%
\definecolor{currentstroke}{rgb}{0.000000,0.000000,0.000000}%
\pgfsetstrokecolor{currentstroke}%
\pgfsetdash{}{0pt}%
\pgfsys@defobject{currentmarker}{\pgfqpoint{0.000000in}{0.000000in}}{\pgfqpoint{0.000000in}{0.020833in}}{%
\pgfpathmoveto{\pgfqpoint{0.000000in}{0.000000in}}%
\pgfpathlineto{\pgfqpoint{0.000000in}{0.020833in}}%
\pgfusepath{stroke,fill}%
}%
\begin{pgfscope}%
\pgfsys@transformshift{5.458054in}{1.775926in}%
\pgfsys@useobject{currentmarker}{}%
\end{pgfscope}%
\end{pgfscope}%
\begin{pgfscope}%
\pgfsetbuttcap%
\pgfsetroundjoin%
\definecolor{currentfill}{rgb}{0.000000,0.000000,0.000000}%
\pgfsetfillcolor{currentfill}%
\pgfsetlinewidth{0.501875pt}%
\definecolor{currentstroke}{rgb}{0.000000,0.000000,0.000000}%
\pgfsetstrokecolor{currentstroke}%
\pgfsetdash{}{0pt}%
\pgfsys@defobject{currentmarker}{\pgfqpoint{0.000000in}{-0.020833in}}{\pgfqpoint{0.000000in}{0.000000in}}{%
\pgfpathmoveto{\pgfqpoint{0.000000in}{0.000000in}}%
\pgfpathlineto{\pgfqpoint{0.000000in}{-0.020833in}}%
\pgfusepath{stroke,fill}%
}%
\begin{pgfscope}%
\pgfsys@transformshift{5.458054in}{3.933333in}%
\pgfsys@useobject{currentmarker}{}%
\end{pgfscope}%
\end{pgfscope}%
\begin{pgfscope}%
\pgfsetbuttcap%
\pgfsetroundjoin%
\definecolor{currentfill}{rgb}{0.000000,0.000000,0.000000}%
\pgfsetfillcolor{currentfill}%
\pgfsetlinewidth{0.501875pt}%
\definecolor{currentstroke}{rgb}{0.000000,0.000000,0.000000}%
\pgfsetstrokecolor{currentstroke}%
\pgfsetdash{}{0pt}%
\pgfsys@defobject{currentmarker}{\pgfqpoint{0.000000in}{0.000000in}}{\pgfqpoint{0.000000in}{0.020833in}}{%
\pgfpathmoveto{\pgfqpoint{0.000000in}{0.000000in}}%
\pgfpathlineto{\pgfqpoint{0.000000in}{0.020833in}}%
\pgfusepath{stroke,fill}%
}%
\begin{pgfscope}%
\pgfsys@transformshift{5.636097in}{1.775926in}%
\pgfsys@useobject{currentmarker}{}%
\end{pgfscope}%
\end{pgfscope}%
\begin{pgfscope}%
\pgfsetbuttcap%
\pgfsetroundjoin%
\definecolor{currentfill}{rgb}{0.000000,0.000000,0.000000}%
\pgfsetfillcolor{currentfill}%
\pgfsetlinewidth{0.501875pt}%
\definecolor{currentstroke}{rgb}{0.000000,0.000000,0.000000}%
\pgfsetstrokecolor{currentstroke}%
\pgfsetdash{}{0pt}%
\pgfsys@defobject{currentmarker}{\pgfqpoint{0.000000in}{-0.020833in}}{\pgfqpoint{0.000000in}{0.000000in}}{%
\pgfpathmoveto{\pgfqpoint{0.000000in}{0.000000in}}%
\pgfpathlineto{\pgfqpoint{0.000000in}{-0.020833in}}%
\pgfusepath{stroke,fill}%
}%
\begin{pgfscope}%
\pgfsys@transformshift{5.636097in}{3.933333in}%
\pgfsys@useobject{currentmarker}{}%
\end{pgfscope}%
\end{pgfscope}%
\begin{pgfscope}%
\pgfsetbuttcap%
\pgfsetroundjoin%
\definecolor{currentfill}{rgb}{0.000000,0.000000,0.000000}%
\pgfsetfillcolor{currentfill}%
\pgfsetlinewidth{0.501875pt}%
\definecolor{currentstroke}{rgb}{0.000000,0.000000,0.000000}%
\pgfsetstrokecolor{currentstroke}%
\pgfsetdash{}{0pt}%
\pgfsys@defobject{currentmarker}{\pgfqpoint{0.000000in}{0.000000in}}{\pgfqpoint{0.000000in}{0.020833in}}{%
\pgfpathmoveto{\pgfqpoint{0.000000in}{0.000000in}}%
\pgfpathlineto{\pgfqpoint{0.000000in}{0.020833in}}%
\pgfusepath{stroke,fill}%
}%
\begin{pgfscope}%
\pgfsys@transformshift{5.814140in}{1.775926in}%
\pgfsys@useobject{currentmarker}{}%
\end{pgfscope}%
\end{pgfscope}%
\begin{pgfscope}%
\pgfsetbuttcap%
\pgfsetroundjoin%
\definecolor{currentfill}{rgb}{0.000000,0.000000,0.000000}%
\pgfsetfillcolor{currentfill}%
\pgfsetlinewidth{0.501875pt}%
\definecolor{currentstroke}{rgb}{0.000000,0.000000,0.000000}%
\pgfsetstrokecolor{currentstroke}%
\pgfsetdash{}{0pt}%
\pgfsys@defobject{currentmarker}{\pgfqpoint{0.000000in}{-0.020833in}}{\pgfqpoint{0.000000in}{0.000000in}}{%
\pgfpathmoveto{\pgfqpoint{0.000000in}{0.000000in}}%
\pgfpathlineto{\pgfqpoint{0.000000in}{-0.020833in}}%
\pgfusepath{stroke,fill}%
}%
\begin{pgfscope}%
\pgfsys@transformshift{5.814140in}{3.933333in}%
\pgfsys@useobject{currentmarker}{}%
\end{pgfscope}%
\end{pgfscope}%
\begin{pgfscope}%
\pgfsetbuttcap%
\pgfsetroundjoin%
\definecolor{currentfill}{rgb}{0.000000,0.000000,0.000000}%
\pgfsetfillcolor{currentfill}%
\pgfsetlinewidth{0.501875pt}%
\definecolor{currentstroke}{rgb}{0.000000,0.000000,0.000000}%
\pgfsetstrokecolor{currentstroke}%
\pgfsetdash{}{0pt}%
\pgfsys@defobject{currentmarker}{\pgfqpoint{0.000000in}{0.000000in}}{\pgfqpoint{0.000000in}{0.020833in}}{%
\pgfpathmoveto{\pgfqpoint{0.000000in}{0.000000in}}%
\pgfpathlineto{\pgfqpoint{0.000000in}{0.020833in}}%
\pgfusepath{stroke,fill}%
}%
\begin{pgfscope}%
\pgfsys@transformshift{5.992183in}{1.775926in}%
\pgfsys@useobject{currentmarker}{}%
\end{pgfscope}%
\end{pgfscope}%
\begin{pgfscope}%
\pgfsetbuttcap%
\pgfsetroundjoin%
\definecolor{currentfill}{rgb}{0.000000,0.000000,0.000000}%
\pgfsetfillcolor{currentfill}%
\pgfsetlinewidth{0.501875pt}%
\definecolor{currentstroke}{rgb}{0.000000,0.000000,0.000000}%
\pgfsetstrokecolor{currentstroke}%
\pgfsetdash{}{0pt}%
\pgfsys@defobject{currentmarker}{\pgfqpoint{0.000000in}{-0.020833in}}{\pgfqpoint{0.000000in}{0.000000in}}{%
\pgfpathmoveto{\pgfqpoint{0.000000in}{0.000000in}}%
\pgfpathlineto{\pgfqpoint{0.000000in}{-0.020833in}}%
\pgfusepath{stroke,fill}%
}%
\begin{pgfscope}%
\pgfsys@transformshift{5.992183in}{3.933333in}%
\pgfsys@useobject{currentmarker}{}%
\end{pgfscope}%
\end{pgfscope}%
\begin{pgfscope}%
\pgfsetbuttcap%
\pgfsetroundjoin%
\definecolor{currentfill}{rgb}{0.000000,0.000000,0.000000}%
\pgfsetfillcolor{currentfill}%
\pgfsetlinewidth{0.501875pt}%
\definecolor{currentstroke}{rgb}{0.000000,0.000000,0.000000}%
\pgfsetstrokecolor{currentstroke}%
\pgfsetdash{}{0pt}%
\pgfsys@defobject{currentmarker}{\pgfqpoint{0.000000in}{0.000000in}}{\pgfqpoint{0.000000in}{0.020833in}}{%
\pgfpathmoveto{\pgfqpoint{0.000000in}{0.000000in}}%
\pgfpathlineto{\pgfqpoint{0.000000in}{0.020833in}}%
\pgfusepath{stroke,fill}%
}%
\begin{pgfscope}%
\pgfsys@transformshift{6.170226in}{1.775926in}%
\pgfsys@useobject{currentmarker}{}%
\end{pgfscope}%
\end{pgfscope}%
\begin{pgfscope}%
\pgfsetbuttcap%
\pgfsetroundjoin%
\definecolor{currentfill}{rgb}{0.000000,0.000000,0.000000}%
\pgfsetfillcolor{currentfill}%
\pgfsetlinewidth{0.501875pt}%
\definecolor{currentstroke}{rgb}{0.000000,0.000000,0.000000}%
\pgfsetstrokecolor{currentstroke}%
\pgfsetdash{}{0pt}%
\pgfsys@defobject{currentmarker}{\pgfqpoint{0.000000in}{-0.020833in}}{\pgfqpoint{0.000000in}{0.000000in}}{%
\pgfpathmoveto{\pgfqpoint{0.000000in}{0.000000in}}%
\pgfpathlineto{\pgfqpoint{0.000000in}{-0.020833in}}%
\pgfusepath{stroke,fill}%
}%
\begin{pgfscope}%
\pgfsys@transformshift{6.170226in}{3.933333in}%
\pgfsys@useobject{currentmarker}{}%
\end{pgfscope}%
\end{pgfscope}%
\begin{pgfscope}%
\pgfsetbuttcap%
\pgfsetroundjoin%
\definecolor{currentfill}{rgb}{0.000000,0.000000,0.000000}%
\pgfsetfillcolor{currentfill}%
\pgfsetlinewidth{0.501875pt}%
\definecolor{currentstroke}{rgb}{0.000000,0.000000,0.000000}%
\pgfsetstrokecolor{currentstroke}%
\pgfsetdash{}{0pt}%
\pgfsys@defobject{currentmarker}{\pgfqpoint{-0.041667in}{-0.000000in}}{\pgfqpoint{0.041667in}{0.000000in}}{%
\pgfpathmoveto{\pgfqpoint{0.041667in}{-0.000000in}}%
\pgfpathlineto{\pgfqpoint{-0.041667in}{0.000000in}}%
\pgfusepath{stroke,fill}%
}%
\begin{pgfscope}%
\pgfsys@transformshift{0.328192in}{0.337654in}%
\pgfsys@useobject{currentmarker}{}%
\end{pgfscope}%
\end{pgfscope}%
\begin{pgfscope}%
\pgfsetbuttcap%
\pgfsetroundjoin%
\definecolor{currentfill}{rgb}{0.000000,0.000000,0.000000}%
\pgfsetfillcolor{currentfill}%
\pgfsetlinewidth{0.501875pt}%
\definecolor{currentstroke}{rgb}{0.000000,0.000000,0.000000}%
\pgfsetstrokecolor{currentstroke}%
\pgfsetdash{}{0pt}%
\pgfsys@defobject{currentmarker}{\pgfqpoint{-0.041667in}{-0.000000in}}{\pgfqpoint{0.041667in}{0.000000in}}{%
\pgfpathmoveto{\pgfqpoint{0.041667in}{-0.000000in}}%
\pgfpathlineto{\pgfqpoint{-0.041667in}{0.000000in}}%
\pgfusepath{stroke,fill}%
}%
\begin{pgfscope}%
\pgfsys@transformshift{6.170226in}{0.337654in}%
\pgfsys@useobject{currentmarker}{}%
\end{pgfscope}%
\end{pgfscope}%
\begin{pgfscope}%
\definecolor{textcolor}{rgb}{0.000000,0.000000,0.000000}%
\pgfsetstrokecolor{textcolor}%
\pgfsetfillcolor{textcolor}%
\pgftext[x=0.084998in, y=0.288956in, left, base]{\color{textcolor}{\rmfamily\fontsize{10.000000}{12.000000}\selectfont\catcode`\^=\active\def^{\ifmmode\sp\else\^{}\fi}\catcode`\%=\active\def%{\%}$-1$}}%
\end{pgfscope}%
\begin{pgfscope}%
\pgfsetbuttcap%
\pgfsetroundjoin%
\definecolor{currentfill}{rgb}{0.000000,0.000000,0.000000}%
\pgfsetfillcolor{currentfill}%
\pgfsetlinewidth{0.501875pt}%
\definecolor{currentstroke}{rgb}{0.000000,0.000000,0.000000}%
\pgfsetstrokecolor{currentstroke}%
\pgfsetdash{}{0pt}%
\pgfsys@defobject{currentmarker}{\pgfqpoint{-0.041667in}{-0.000000in}}{\pgfqpoint{0.041667in}{0.000000in}}{%
\pgfpathmoveto{\pgfqpoint{0.041667in}{-0.000000in}}%
\pgfpathlineto{\pgfqpoint{-0.041667in}{0.000000in}}%
\pgfusepath{stroke,fill}%
}%
\begin{pgfscope}%
\pgfsys@transformshift{0.328192in}{3.214198in}%
\pgfsys@useobject{currentmarker}{}%
\end{pgfscope}%
\end{pgfscope}%
\begin{pgfscope}%
\pgfsetbuttcap%
\pgfsetroundjoin%
\definecolor{currentfill}{rgb}{0.000000,0.000000,0.000000}%
\pgfsetfillcolor{currentfill}%
\pgfsetlinewidth{0.501875pt}%
\definecolor{currentstroke}{rgb}{0.000000,0.000000,0.000000}%
\pgfsetstrokecolor{currentstroke}%
\pgfsetdash{}{0pt}%
\pgfsys@defobject{currentmarker}{\pgfqpoint{-0.041667in}{-0.000000in}}{\pgfqpoint{0.041667in}{0.000000in}}{%
\pgfpathmoveto{\pgfqpoint{0.041667in}{-0.000000in}}%
\pgfpathlineto{\pgfqpoint{-0.041667in}{0.000000in}}%
\pgfusepath{stroke,fill}%
}%
\begin{pgfscope}%
\pgfsys@transformshift{6.170226in}{3.214198in}%
\pgfsys@useobject{currentmarker}{}%
\end{pgfscope}%
\end{pgfscope}%
\begin{pgfscope}%
\definecolor{textcolor}{rgb}{0.000000,0.000000,0.000000}%
\pgfsetstrokecolor{textcolor}%
\pgfsetfillcolor{textcolor}%
\pgftext[x=0.173331in, y=3.165500in, left, base]{\color{textcolor}{\rmfamily\fontsize{10.000000}{12.000000}\selectfont\catcode`\^=\active\def^{\ifmmode\sp\else\^{}\fi}\catcode`\%=\active\def%{\%}$1$}}%
\end{pgfscope}%
\begin{pgfscope}%
\pgfsetbuttcap%
\pgfsetroundjoin%
\definecolor{currentfill}{rgb}{0.000000,0.000000,0.000000}%
\pgfsetfillcolor{currentfill}%
\pgfsetlinewidth{0.501875pt}%
\definecolor{currentstroke}{rgb}{0.000000,0.000000,0.000000}%
\pgfsetstrokecolor{currentstroke}%
\pgfsetdash{}{0pt}%
\pgfsys@defobject{currentmarker}{\pgfqpoint{0.000000in}{0.000000in}}{\pgfqpoint{0.020833in}{0.000000in}}{%
\pgfpathmoveto{\pgfqpoint{0.000000in}{0.000000in}}%
\pgfpathlineto{\pgfqpoint{0.020833in}{0.000000in}}%
\pgfusepath{stroke,fill}%
}%
\begin{pgfscope}%
\pgfsys@transformshift{0.328192in}{1.056790in}%
\pgfsys@useobject{currentmarker}{}%
\end{pgfscope}%
\end{pgfscope}%
\begin{pgfscope}%
\pgfsetbuttcap%
\pgfsetroundjoin%
\definecolor{currentfill}{rgb}{0.000000,0.000000,0.000000}%
\pgfsetfillcolor{currentfill}%
\pgfsetlinewidth{0.501875pt}%
\definecolor{currentstroke}{rgb}{0.000000,0.000000,0.000000}%
\pgfsetstrokecolor{currentstroke}%
\pgfsetdash{}{0pt}%
\pgfsys@defobject{currentmarker}{\pgfqpoint{-0.020833in}{0.000000in}}{\pgfqpoint{-0.000000in}{0.000000in}}{%
\pgfpathmoveto{\pgfqpoint{-0.000000in}{0.000000in}}%
\pgfpathlineto{\pgfqpoint{-0.020833in}{0.000000in}}%
\pgfusepath{stroke,fill}%
}%
\begin{pgfscope}%
\pgfsys@transformshift{6.170226in}{1.056790in}%
\pgfsys@useobject{currentmarker}{}%
\end{pgfscope}%
\end{pgfscope}%
\begin{pgfscope}%
\pgfsetbuttcap%
\pgfsetroundjoin%
\definecolor{currentfill}{rgb}{0.000000,0.000000,0.000000}%
\pgfsetfillcolor{currentfill}%
\pgfsetlinewidth{0.501875pt}%
\definecolor{currentstroke}{rgb}{0.000000,0.000000,0.000000}%
\pgfsetstrokecolor{currentstroke}%
\pgfsetdash{}{0pt}%
\pgfsys@defobject{currentmarker}{\pgfqpoint{0.000000in}{0.000000in}}{\pgfqpoint{0.020833in}{0.000000in}}{%
\pgfpathmoveto{\pgfqpoint{0.000000in}{0.000000in}}%
\pgfpathlineto{\pgfqpoint{0.020833in}{0.000000in}}%
\pgfusepath{stroke,fill}%
}%
\begin{pgfscope}%
\pgfsys@transformshift{0.328192in}{1.775926in}%
\pgfsys@useobject{currentmarker}{}%
\end{pgfscope}%
\end{pgfscope}%
\begin{pgfscope}%
\pgfsetbuttcap%
\pgfsetroundjoin%
\definecolor{currentfill}{rgb}{0.000000,0.000000,0.000000}%
\pgfsetfillcolor{currentfill}%
\pgfsetlinewidth{0.501875pt}%
\definecolor{currentstroke}{rgb}{0.000000,0.000000,0.000000}%
\pgfsetstrokecolor{currentstroke}%
\pgfsetdash{}{0pt}%
\pgfsys@defobject{currentmarker}{\pgfqpoint{-0.020833in}{0.000000in}}{\pgfqpoint{-0.000000in}{0.000000in}}{%
\pgfpathmoveto{\pgfqpoint{-0.000000in}{0.000000in}}%
\pgfpathlineto{\pgfqpoint{-0.020833in}{0.000000in}}%
\pgfusepath{stroke,fill}%
}%
\begin{pgfscope}%
\pgfsys@transformshift{6.170226in}{1.775926in}%
\pgfsys@useobject{currentmarker}{}%
\end{pgfscope}%
\end{pgfscope}%
\begin{pgfscope}%
\pgfsetbuttcap%
\pgfsetroundjoin%
\definecolor{currentfill}{rgb}{0.000000,0.000000,0.000000}%
\pgfsetfillcolor{currentfill}%
\pgfsetlinewidth{0.501875pt}%
\definecolor{currentstroke}{rgb}{0.000000,0.000000,0.000000}%
\pgfsetstrokecolor{currentstroke}%
\pgfsetdash{}{0pt}%
\pgfsys@defobject{currentmarker}{\pgfqpoint{0.000000in}{0.000000in}}{\pgfqpoint{0.020833in}{0.000000in}}{%
\pgfpathmoveto{\pgfqpoint{0.000000in}{0.000000in}}%
\pgfpathlineto{\pgfqpoint{0.020833in}{0.000000in}}%
\pgfusepath{stroke,fill}%
}%
\begin{pgfscope}%
\pgfsys@transformshift{0.328192in}{2.495062in}%
\pgfsys@useobject{currentmarker}{}%
\end{pgfscope}%
\end{pgfscope}%
\begin{pgfscope}%
\pgfsetbuttcap%
\pgfsetroundjoin%
\definecolor{currentfill}{rgb}{0.000000,0.000000,0.000000}%
\pgfsetfillcolor{currentfill}%
\pgfsetlinewidth{0.501875pt}%
\definecolor{currentstroke}{rgb}{0.000000,0.000000,0.000000}%
\pgfsetstrokecolor{currentstroke}%
\pgfsetdash{}{0pt}%
\pgfsys@defobject{currentmarker}{\pgfqpoint{-0.020833in}{0.000000in}}{\pgfqpoint{-0.000000in}{0.000000in}}{%
\pgfpathmoveto{\pgfqpoint{-0.000000in}{0.000000in}}%
\pgfpathlineto{\pgfqpoint{-0.020833in}{0.000000in}}%
\pgfusepath{stroke,fill}%
}%
\begin{pgfscope}%
\pgfsys@transformshift{6.170226in}{2.495062in}%
\pgfsys@useobject{currentmarker}{}%
\end{pgfscope}%
\end{pgfscope}%
\begin{pgfscope}%
\pgfsetbuttcap%
\pgfsetroundjoin%
\definecolor{currentfill}{rgb}{0.000000,0.000000,0.000000}%
\pgfsetfillcolor{currentfill}%
\pgfsetlinewidth{0.501875pt}%
\definecolor{currentstroke}{rgb}{0.000000,0.000000,0.000000}%
\pgfsetstrokecolor{currentstroke}%
\pgfsetdash{}{0pt}%
\pgfsys@defobject{currentmarker}{\pgfqpoint{0.000000in}{0.000000in}}{\pgfqpoint{0.020833in}{0.000000in}}{%
\pgfpathmoveto{\pgfqpoint{0.000000in}{0.000000in}}%
\pgfpathlineto{\pgfqpoint{0.020833in}{0.000000in}}%
\pgfusepath{stroke,fill}%
}%
\begin{pgfscope}%
\pgfsys@transformshift{0.328192in}{3.933333in}%
\pgfsys@useobject{currentmarker}{}%
\end{pgfscope}%
\end{pgfscope}%
\begin{pgfscope}%
\pgfsetbuttcap%
\pgfsetroundjoin%
\definecolor{currentfill}{rgb}{0.000000,0.000000,0.000000}%
\pgfsetfillcolor{currentfill}%
\pgfsetlinewidth{0.501875pt}%
\definecolor{currentstroke}{rgb}{0.000000,0.000000,0.000000}%
\pgfsetstrokecolor{currentstroke}%
\pgfsetdash{}{0pt}%
\pgfsys@defobject{currentmarker}{\pgfqpoint{-0.020833in}{0.000000in}}{\pgfqpoint{-0.000000in}{0.000000in}}{%
\pgfpathmoveto{\pgfqpoint{-0.000000in}{0.000000in}}%
\pgfpathlineto{\pgfqpoint{-0.020833in}{0.000000in}}%
\pgfusepath{stroke,fill}%
}%
\begin{pgfscope}%
\pgfsys@transformshift{6.170226in}{3.933333in}%
\pgfsys@useobject{currentmarker}{}%
\end{pgfscope}%
\end{pgfscope}%
\begin{pgfscope}%
\pgfpathrectangle{\pgfqpoint{0.050000in}{0.050000in}}{\pgfqpoint{6.120226in}{3.883333in}}%
\pgfusepath{clip}%
\pgfsetrectcap%
\pgfsetroundjoin%
\pgfsetlinewidth{1.003750pt}%
\definecolor{currentstroke}{rgb}{0.000000,0.000000,0.000000}%
\pgfsetstrokecolor{currentstroke}%
\pgfsetdash{}{0pt}%
\pgfpathmoveto{\pgfqpoint{0.328192in}{1.775926in}}%
\pgfpathlineto{\pgfqpoint{0.329583in}{1.869091in}}%
\pgfpathlineto{\pgfqpoint{0.333756in}{1.643159in}}%
\pgfpathlineto{\pgfqpoint{0.336538in}{1.673075in}}%
\pgfpathlineto{\pgfqpoint{0.337929in}{1.651870in}}%
\pgfpathlineto{\pgfqpoint{0.340711in}{1.720720in}}%
\pgfpathlineto{\pgfqpoint{0.342102in}{1.717390in}}%
\pgfpathlineto{\pgfqpoint{0.344884in}{1.769543in}}%
\pgfpathlineto{\pgfqpoint{0.346275in}{1.701911in}}%
\pgfpathlineto{\pgfqpoint{0.349057in}{1.863470in}}%
\pgfpathlineto{\pgfqpoint{0.350447in}{1.840295in}}%
\pgfpathlineto{\pgfqpoint{0.351838in}{1.872018in}}%
\pgfpathlineto{\pgfqpoint{0.354620in}{1.759381in}}%
\pgfpathlineto{\pgfqpoint{0.356011in}{1.808356in}}%
\pgfpathlineto{\pgfqpoint{0.357402in}{1.789153in}}%
\pgfpathlineto{\pgfqpoint{0.358793in}{1.808522in}}%
\pgfpathlineto{\pgfqpoint{0.361575in}{1.732913in}}%
\pgfpathlineto{\pgfqpoint{0.362966in}{1.761638in}}%
\pgfpathlineto{\pgfqpoint{0.365748in}{1.662019in}}%
\pgfpathlineto{\pgfqpoint{0.367139in}{1.711833in}}%
\pgfpathlineto{\pgfqpoint{0.368530in}{1.704232in}}%
\pgfpathlineto{\pgfqpoint{0.369921in}{1.663477in}}%
\pgfpathlineto{\pgfqpoint{0.372703in}{1.687566in}}%
\pgfpathlineto{\pgfqpoint{0.374094in}{1.674132in}}%
\pgfpathlineto{\pgfqpoint{0.375485in}{1.614007in}}%
\pgfpathlineto{\pgfqpoint{0.376876in}{1.638192in}}%
\pgfpathlineto{\pgfqpoint{0.378267in}{1.628556in}}%
\pgfpathlineto{\pgfqpoint{0.379658in}{1.681507in}}%
\pgfpathlineto{\pgfqpoint{0.381049in}{1.676216in}}%
\pgfpathlineto{\pgfqpoint{0.382440in}{1.689203in}}%
\pgfpathlineto{\pgfqpoint{0.383831in}{1.743552in}}%
\pgfpathlineto{\pgfqpoint{0.385221in}{1.733862in}}%
\pgfpathlineto{\pgfqpoint{0.388003in}{1.616919in}}%
\pgfpathlineto{\pgfqpoint{0.390785in}{1.711013in}}%
\pgfpathlineto{\pgfqpoint{0.394958in}{1.505986in}}%
\pgfpathlineto{\pgfqpoint{0.396349in}{1.531276in}}%
\pgfpathlineto{\pgfqpoint{0.397740in}{1.526505in}}%
\pgfpathlineto{\pgfqpoint{0.400522in}{1.560218in}}%
\pgfpathlineto{\pgfqpoint{0.403304in}{1.419273in}}%
\pgfpathlineto{\pgfqpoint{0.404695in}{1.417287in}}%
\pgfpathlineto{\pgfqpoint{0.407477in}{1.293565in}}%
\pgfpathlineto{\pgfqpoint{0.410259in}{1.335118in}}%
\pgfpathlineto{\pgfqpoint{0.411650in}{1.328908in}}%
\pgfpathlineto{\pgfqpoint{0.413041in}{1.342557in}}%
\pgfpathlineto{\pgfqpoint{0.414432in}{1.340615in}}%
\pgfpathlineto{\pgfqpoint{0.418605in}{1.242877in}}%
\pgfpathlineto{\pgfqpoint{0.419995in}{1.248690in}}%
\pgfpathlineto{\pgfqpoint{0.421386in}{1.271679in}}%
\pgfpathlineto{\pgfqpoint{0.422777in}{1.344988in}}%
\pgfpathlineto{\pgfqpoint{0.424168in}{1.310635in}}%
\pgfpathlineto{\pgfqpoint{0.426950in}{1.369341in}}%
\pgfpathlineto{\pgfqpoint{0.428341in}{1.379244in}}%
\pgfpathlineto{\pgfqpoint{0.429732in}{1.426959in}}%
\pgfpathlineto{\pgfqpoint{0.433905in}{1.288776in}}%
\pgfpathlineto{\pgfqpoint{0.436687in}{1.326966in}}%
\pgfpathlineto{\pgfqpoint{0.438078in}{1.372691in}}%
\pgfpathlineto{\pgfqpoint{0.439469in}{1.229609in}}%
\pgfpathlineto{\pgfqpoint{0.440860in}{1.262489in}}%
\pgfpathlineto{\pgfqpoint{0.442251in}{1.238314in}}%
\pgfpathlineto{\pgfqpoint{0.443642in}{1.242459in}}%
\pgfpathlineto{\pgfqpoint{0.446424in}{1.308235in}}%
\pgfpathlineto{\pgfqpoint{0.447815in}{1.309908in}}%
\pgfpathlineto{\pgfqpoint{0.449206in}{1.307634in}}%
\pgfpathlineto{\pgfqpoint{0.450597in}{1.326909in}}%
\pgfpathlineto{\pgfqpoint{0.451988in}{1.326977in}}%
\pgfpathlineto{\pgfqpoint{0.453379in}{1.221421in}}%
\pgfpathlineto{\pgfqpoint{0.457551in}{1.272409in}}%
\pgfpathlineto{\pgfqpoint{0.458942in}{1.265451in}}%
\pgfpathlineto{\pgfqpoint{0.460333in}{1.350988in}}%
\pgfpathlineto{\pgfqpoint{0.461724in}{1.305518in}}%
\pgfpathlineto{\pgfqpoint{0.463115in}{1.355304in}}%
\pgfpathlineto{\pgfqpoint{0.465897in}{1.267662in}}%
\pgfpathlineto{\pgfqpoint{0.467288in}{1.316726in}}%
\pgfpathlineto{\pgfqpoint{0.468679in}{1.318681in}}%
\pgfpathlineto{\pgfqpoint{0.470070in}{1.298596in}}%
\pgfpathlineto{\pgfqpoint{0.471461in}{1.295937in}}%
\pgfpathlineto{\pgfqpoint{0.472852in}{1.329128in}}%
\pgfpathlineto{\pgfqpoint{0.474243in}{1.326314in}}%
\pgfpathlineto{\pgfqpoint{0.477025in}{1.303952in}}%
\pgfpathlineto{\pgfqpoint{0.478416in}{1.235416in}}%
\pgfpathlineto{\pgfqpoint{0.479807in}{1.266617in}}%
\pgfpathlineto{\pgfqpoint{0.481198in}{1.256946in}}%
\pgfpathlineto{\pgfqpoint{0.482589in}{1.294073in}}%
\pgfpathlineto{\pgfqpoint{0.483980in}{1.262776in}}%
\pgfpathlineto{\pgfqpoint{0.485371in}{1.280993in}}%
\pgfpathlineto{\pgfqpoint{0.486762in}{1.238139in}}%
\pgfpathlineto{\pgfqpoint{0.488153in}{1.254722in}}%
\pgfpathlineto{\pgfqpoint{0.489544in}{1.245921in}}%
\pgfpathlineto{\pgfqpoint{0.492325in}{1.211086in}}%
\pgfpathlineto{\pgfqpoint{0.496498in}{1.375515in}}%
\pgfpathlineto{\pgfqpoint{0.497889in}{1.405220in}}%
\pgfpathlineto{\pgfqpoint{0.500671in}{1.351203in}}%
\pgfpathlineto{\pgfqpoint{0.502062in}{1.329056in}}%
\pgfpathlineto{\pgfqpoint{0.503453in}{1.420623in}}%
\pgfpathlineto{\pgfqpoint{0.504844in}{1.387161in}}%
\pgfpathlineto{\pgfqpoint{0.506235in}{1.396822in}}%
\pgfpathlineto{\pgfqpoint{0.509017in}{1.320497in}}%
\pgfpathlineto{\pgfqpoint{0.511799in}{1.477203in}}%
\pgfpathlineto{\pgfqpoint{0.513190in}{1.441717in}}%
\pgfpathlineto{\pgfqpoint{0.514581in}{1.459483in}}%
\pgfpathlineto{\pgfqpoint{0.515972in}{1.459870in}}%
\pgfpathlineto{\pgfqpoint{0.517363in}{1.451080in}}%
\pgfpathlineto{\pgfqpoint{0.518754in}{1.414727in}}%
\pgfpathlineto{\pgfqpoint{0.520145in}{1.336756in}}%
\pgfpathlineto{\pgfqpoint{0.521536in}{1.348148in}}%
\pgfpathlineto{\pgfqpoint{0.522927in}{1.334226in}}%
\pgfpathlineto{\pgfqpoint{0.524318in}{1.307143in}}%
\pgfpathlineto{\pgfqpoint{0.525708in}{1.231678in}}%
\pgfpathlineto{\pgfqpoint{0.527099in}{1.241632in}}%
\pgfpathlineto{\pgfqpoint{0.528490in}{1.348228in}}%
\pgfpathlineto{\pgfqpoint{0.529881in}{1.354141in}}%
\pgfpathlineto{\pgfqpoint{0.531272in}{1.352831in}}%
\pgfpathlineto{\pgfqpoint{0.532663in}{1.370799in}}%
\pgfpathlineto{\pgfqpoint{0.534054in}{1.437893in}}%
\pgfpathlineto{\pgfqpoint{0.535445in}{1.430018in}}%
\pgfpathlineto{\pgfqpoint{0.536836in}{1.460698in}}%
\pgfpathlineto{\pgfqpoint{0.538227in}{1.543167in}}%
\pgfpathlineto{\pgfqpoint{0.539618in}{1.529427in}}%
\pgfpathlineto{\pgfqpoint{0.541009in}{1.415421in}}%
\pgfpathlineto{\pgfqpoint{0.543791in}{1.557014in}}%
\pgfpathlineto{\pgfqpoint{0.546573in}{1.490450in}}%
\pgfpathlineto{\pgfqpoint{0.547964in}{1.500718in}}%
\pgfpathlineto{\pgfqpoint{0.549355in}{1.499004in}}%
\pgfpathlineto{\pgfqpoint{0.550746in}{1.503150in}}%
\pgfpathlineto{\pgfqpoint{0.552137in}{1.516796in}}%
\pgfpathlineto{\pgfqpoint{0.553528in}{1.493443in}}%
\pgfpathlineto{\pgfqpoint{0.556310in}{1.588975in}}%
\pgfpathlineto{\pgfqpoint{0.557701in}{1.604141in}}%
\pgfpathlineto{\pgfqpoint{0.559092in}{1.591684in}}%
\pgfpathlineto{\pgfqpoint{0.560483in}{1.625609in}}%
\pgfpathlineto{\pgfqpoint{0.563264in}{1.747249in}}%
\pgfpathlineto{\pgfqpoint{0.564655in}{1.671875in}}%
\pgfpathlineto{\pgfqpoint{0.566046in}{1.674981in}}%
\pgfpathlineto{\pgfqpoint{0.570219in}{1.588596in}}%
\pgfpathlineto{\pgfqpoint{0.571610in}{1.643468in}}%
\pgfpathlineto{\pgfqpoint{0.574392in}{1.611044in}}%
\pgfpathlineto{\pgfqpoint{0.577174in}{1.722267in}}%
\pgfpathlineto{\pgfqpoint{0.578565in}{1.708147in}}%
\pgfpathlineto{\pgfqpoint{0.579956in}{1.635528in}}%
\pgfpathlineto{\pgfqpoint{0.581347in}{1.693483in}}%
\pgfpathlineto{\pgfqpoint{0.582738in}{1.633297in}}%
\pgfpathlineto{\pgfqpoint{0.586911in}{1.716253in}}%
\pgfpathlineto{\pgfqpoint{0.588302in}{1.769563in}}%
\pgfpathlineto{\pgfqpoint{0.591084in}{1.737266in}}%
\pgfpathlineto{\pgfqpoint{0.592475in}{1.804516in}}%
\pgfpathlineto{\pgfqpoint{0.593866in}{1.791478in}}%
\pgfpathlineto{\pgfqpoint{0.596647in}{1.659218in}}%
\pgfpathlineto{\pgfqpoint{0.602211in}{1.860631in}}%
\pgfpathlineto{\pgfqpoint{0.603602in}{1.850728in}}%
\pgfpathlineto{\pgfqpoint{0.607775in}{1.665162in}}%
\pgfpathlineto{\pgfqpoint{0.609166in}{1.763775in}}%
\pgfpathlineto{\pgfqpoint{0.610557in}{1.700067in}}%
\pgfpathlineto{\pgfqpoint{0.613339in}{1.752222in}}%
\pgfpathlineto{\pgfqpoint{0.614730in}{1.743430in}}%
\pgfpathlineto{\pgfqpoint{0.616121in}{1.750162in}}%
\pgfpathlineto{\pgfqpoint{0.617512in}{1.847006in}}%
\pgfpathlineto{\pgfqpoint{0.620294in}{1.766739in}}%
\pgfpathlineto{\pgfqpoint{0.628640in}{2.108006in}}%
\pgfpathlineto{\pgfqpoint{0.630031in}{2.127320in}}%
\pgfpathlineto{\pgfqpoint{0.632812in}{2.219276in}}%
\pgfpathlineto{\pgfqpoint{0.636985in}{2.104339in}}%
\pgfpathlineto{\pgfqpoint{0.638376in}{2.085694in}}%
\pgfpathlineto{\pgfqpoint{0.639767in}{2.040085in}}%
\pgfpathlineto{\pgfqpoint{0.641158in}{2.081074in}}%
\pgfpathlineto{\pgfqpoint{0.642549in}{2.082769in}}%
\pgfpathlineto{\pgfqpoint{0.643940in}{2.094011in}}%
\pgfpathlineto{\pgfqpoint{0.645331in}{2.086740in}}%
\pgfpathlineto{\pgfqpoint{0.648113in}{2.212580in}}%
\pgfpathlineto{\pgfqpoint{0.649504in}{2.153596in}}%
\pgfpathlineto{\pgfqpoint{0.652286in}{2.220770in}}%
\pgfpathlineto{\pgfqpoint{0.653677in}{2.253952in}}%
\pgfpathlineto{\pgfqpoint{0.655068in}{2.238202in}}%
\pgfpathlineto{\pgfqpoint{0.656459in}{2.333593in}}%
\pgfpathlineto{\pgfqpoint{0.657850in}{2.336678in}}%
\pgfpathlineto{\pgfqpoint{0.660632in}{2.416463in}}%
\pgfpathlineto{\pgfqpoint{0.662023in}{2.336442in}}%
\pgfpathlineto{\pgfqpoint{0.663414in}{2.394281in}}%
\pgfpathlineto{\pgfqpoint{0.664805in}{2.381016in}}%
\pgfpathlineto{\pgfqpoint{0.666195in}{2.316594in}}%
\pgfpathlineto{\pgfqpoint{0.667586in}{2.430526in}}%
\pgfpathlineto{\pgfqpoint{0.668977in}{2.408577in}}%
\pgfpathlineto{\pgfqpoint{0.670368in}{2.418674in}}%
\pgfpathlineto{\pgfqpoint{0.671759in}{2.414559in}}%
\pgfpathlineto{\pgfqpoint{0.673150in}{2.447743in}}%
\pgfpathlineto{\pgfqpoint{0.674541in}{2.418888in}}%
\pgfpathlineto{\pgfqpoint{0.677323in}{2.333369in}}%
\pgfpathlineto{\pgfqpoint{0.678714in}{2.355650in}}%
\pgfpathlineto{\pgfqpoint{0.680105in}{2.419616in}}%
\pgfpathlineto{\pgfqpoint{0.681496in}{2.402523in}}%
\pgfpathlineto{\pgfqpoint{0.682887in}{2.451590in}}%
\pgfpathlineto{\pgfqpoint{0.684278in}{2.407288in}}%
\pgfpathlineto{\pgfqpoint{0.685669in}{2.419568in}}%
\pgfpathlineto{\pgfqpoint{0.687060in}{2.421449in}}%
\pgfpathlineto{\pgfqpoint{0.688451in}{2.498773in}}%
\pgfpathlineto{\pgfqpoint{0.689842in}{2.511739in}}%
\pgfpathlineto{\pgfqpoint{0.694015in}{2.444383in}}%
\pgfpathlineto{\pgfqpoint{0.695406in}{2.446265in}}%
\pgfpathlineto{\pgfqpoint{0.698188in}{2.299985in}}%
\pgfpathlineto{\pgfqpoint{0.699579in}{2.305183in}}%
\pgfpathlineto{\pgfqpoint{0.702360in}{2.234982in}}%
\pgfpathlineto{\pgfqpoint{0.703751in}{2.242127in}}%
\pgfpathlineto{\pgfqpoint{0.705142in}{2.288813in}}%
\pgfpathlineto{\pgfqpoint{0.706533in}{2.236274in}}%
\pgfpathlineto{\pgfqpoint{0.707924in}{2.306147in}}%
\pgfpathlineto{\pgfqpoint{0.709315in}{2.291760in}}%
\pgfpathlineto{\pgfqpoint{0.712097in}{2.369770in}}%
\pgfpathlineto{\pgfqpoint{0.713488in}{2.308811in}}%
\pgfpathlineto{\pgfqpoint{0.714879in}{2.333620in}}%
\pgfpathlineto{\pgfqpoint{0.716270in}{2.324213in}}%
\pgfpathlineto{\pgfqpoint{0.719052in}{2.249932in}}%
\pgfpathlineto{\pgfqpoint{0.720443in}{2.283997in}}%
\pgfpathlineto{\pgfqpoint{0.721834in}{2.261686in}}%
\pgfpathlineto{\pgfqpoint{0.723225in}{2.265653in}}%
\pgfpathlineto{\pgfqpoint{0.726007in}{2.231585in}}%
\pgfpathlineto{\pgfqpoint{0.727398in}{2.212590in}}%
\pgfpathlineto{\pgfqpoint{0.728789in}{2.249686in}}%
\pgfpathlineto{\pgfqpoint{0.730180in}{2.179873in}}%
\pgfpathlineto{\pgfqpoint{0.731571in}{2.225833in}}%
\pgfpathlineto{\pgfqpoint{0.732962in}{2.231423in}}%
\pgfpathlineto{\pgfqpoint{0.734353in}{2.346836in}}%
\pgfpathlineto{\pgfqpoint{0.735744in}{2.291993in}}%
\pgfpathlineto{\pgfqpoint{0.737134in}{2.387362in}}%
\pgfpathlineto{\pgfqpoint{0.738525in}{2.361662in}}%
\pgfpathlineto{\pgfqpoint{0.741307in}{2.389274in}}%
\pgfpathlineto{\pgfqpoint{0.744089in}{2.450591in}}%
\pgfpathlineto{\pgfqpoint{0.745480in}{2.435277in}}%
\pgfpathlineto{\pgfqpoint{0.746871in}{2.457177in}}%
\pgfpathlineto{\pgfqpoint{0.748262in}{2.530570in}}%
\pgfpathlineto{\pgfqpoint{0.749653in}{2.535081in}}%
\pgfpathlineto{\pgfqpoint{0.752435in}{2.622877in}}%
\pgfpathlineto{\pgfqpoint{0.753826in}{2.558825in}}%
\pgfpathlineto{\pgfqpoint{0.755217in}{2.570145in}}%
\pgfpathlineto{\pgfqpoint{0.756608in}{2.535278in}}%
\pgfpathlineto{\pgfqpoint{0.757999in}{2.535165in}}%
\pgfpathlineto{\pgfqpoint{0.760781in}{2.451510in}}%
\pgfpathlineto{\pgfqpoint{0.762172in}{2.466611in}}%
\pgfpathlineto{\pgfqpoint{0.764954in}{2.374393in}}%
\pgfpathlineto{\pgfqpoint{0.770518in}{2.499375in}}%
\pgfpathlineto{\pgfqpoint{0.773299in}{2.443167in}}%
\pgfpathlineto{\pgfqpoint{0.774690in}{2.386358in}}%
\pgfpathlineto{\pgfqpoint{0.777472in}{2.427767in}}%
\pgfpathlineto{\pgfqpoint{0.778863in}{2.468770in}}%
\pgfpathlineto{\pgfqpoint{0.780254in}{2.469447in}}%
\pgfpathlineto{\pgfqpoint{0.781645in}{2.459681in}}%
\pgfpathlineto{\pgfqpoint{0.783036in}{2.377256in}}%
\pgfpathlineto{\pgfqpoint{0.785818in}{2.455801in}}%
\pgfpathlineto{\pgfqpoint{0.788600in}{2.382302in}}%
\pgfpathlineto{\pgfqpoint{0.789991in}{2.410138in}}%
\pgfpathlineto{\pgfqpoint{0.791382in}{2.382131in}}%
\pgfpathlineto{\pgfqpoint{0.792773in}{2.444471in}}%
\pgfpathlineto{\pgfqpoint{0.794164in}{2.442896in}}%
\pgfpathlineto{\pgfqpoint{0.795555in}{2.438403in}}%
\pgfpathlineto{\pgfqpoint{0.796946in}{2.449460in}}%
\pgfpathlineto{\pgfqpoint{0.798337in}{2.488677in}}%
\pgfpathlineto{\pgfqpoint{0.801119in}{2.406809in}}%
\pgfpathlineto{\pgfqpoint{0.802510in}{2.414457in}}%
\pgfpathlineto{\pgfqpoint{0.803901in}{2.394076in}}%
\pgfpathlineto{\pgfqpoint{0.805292in}{2.336269in}}%
\pgfpathlineto{\pgfqpoint{0.808073in}{2.377975in}}%
\pgfpathlineto{\pgfqpoint{0.809464in}{2.345051in}}%
\pgfpathlineto{\pgfqpoint{0.810855in}{2.346192in}}%
\pgfpathlineto{\pgfqpoint{0.812246in}{2.282036in}}%
\pgfpathlineto{\pgfqpoint{0.813637in}{2.330181in}}%
\pgfpathlineto{\pgfqpoint{0.815028in}{2.337234in}}%
\pgfpathlineto{\pgfqpoint{0.816419in}{2.378356in}}%
\pgfpathlineto{\pgfqpoint{0.817810in}{2.336452in}}%
\pgfpathlineto{\pgfqpoint{0.819201in}{2.347317in}}%
\pgfpathlineto{\pgfqpoint{0.820592in}{2.417797in}}%
\pgfpathlineto{\pgfqpoint{0.824765in}{2.349235in}}%
\pgfpathlineto{\pgfqpoint{0.826156in}{2.388760in}}%
\pgfpathlineto{\pgfqpoint{0.828938in}{2.307353in}}%
\pgfpathlineto{\pgfqpoint{0.830329in}{2.340997in}}%
\pgfpathlineto{\pgfqpoint{0.831720in}{2.315624in}}%
\pgfpathlineto{\pgfqpoint{0.833111in}{2.391130in}}%
\pgfpathlineto{\pgfqpoint{0.834502in}{2.383055in}}%
\pgfpathlineto{\pgfqpoint{0.835893in}{2.406699in}}%
\pgfpathlineto{\pgfqpoint{0.837284in}{2.468678in}}%
\pgfpathlineto{\pgfqpoint{0.840066in}{2.428046in}}%
\pgfpathlineto{\pgfqpoint{0.841457in}{2.428027in}}%
\pgfpathlineto{\pgfqpoint{0.842847in}{2.489408in}}%
\pgfpathlineto{\pgfqpoint{0.844238in}{2.461951in}}%
\pgfpathlineto{\pgfqpoint{0.845629in}{2.489052in}}%
\pgfpathlineto{\pgfqpoint{0.849802in}{2.380839in}}%
\pgfpathlineto{\pgfqpoint{0.851193in}{2.376401in}}%
\pgfpathlineto{\pgfqpoint{0.852584in}{2.352625in}}%
\pgfpathlineto{\pgfqpoint{0.855366in}{2.431688in}}%
\pgfpathlineto{\pgfqpoint{0.859539in}{2.548578in}}%
\pgfpathlineto{\pgfqpoint{0.860930in}{2.562070in}}%
\pgfpathlineto{\pgfqpoint{0.862321in}{2.560747in}}%
\pgfpathlineto{\pgfqpoint{0.863712in}{2.551167in}}%
\pgfpathlineto{\pgfqpoint{0.865103in}{2.582637in}}%
\pgfpathlineto{\pgfqpoint{0.866494in}{2.578430in}}%
\pgfpathlineto{\pgfqpoint{0.867885in}{2.530247in}}%
\pgfpathlineto{\pgfqpoint{0.869276in}{2.532846in}}%
\pgfpathlineto{\pgfqpoint{0.872058in}{2.464832in}}%
\pgfpathlineto{\pgfqpoint{0.873449in}{2.482529in}}%
\pgfpathlineto{\pgfqpoint{0.874840in}{2.447635in}}%
\pgfpathlineto{\pgfqpoint{0.876231in}{2.472083in}}%
\pgfpathlineto{\pgfqpoint{0.877621in}{2.440602in}}%
\pgfpathlineto{\pgfqpoint{0.880403in}{2.582106in}}%
\pgfpathlineto{\pgfqpoint{0.881794in}{2.576088in}}%
\pgfpathlineto{\pgfqpoint{0.884576in}{2.594847in}}%
\pgfpathlineto{\pgfqpoint{0.885967in}{2.588352in}}%
\pgfpathlineto{\pgfqpoint{0.887358in}{2.639151in}}%
\pgfpathlineto{\pgfqpoint{0.890140in}{2.559031in}}%
\pgfpathlineto{\pgfqpoint{0.891531in}{2.640619in}}%
\pgfpathlineto{\pgfqpoint{0.894313in}{2.588945in}}%
\pgfpathlineto{\pgfqpoint{0.897095in}{2.668883in}}%
\pgfpathlineto{\pgfqpoint{0.898486in}{2.664748in}}%
\pgfpathlineto{\pgfqpoint{0.899877in}{2.613893in}}%
\pgfpathlineto{\pgfqpoint{0.902659in}{2.743174in}}%
\pgfpathlineto{\pgfqpoint{0.904050in}{2.729245in}}%
\pgfpathlineto{\pgfqpoint{0.905441in}{2.644059in}}%
\pgfpathlineto{\pgfqpoint{0.906832in}{2.668151in}}%
\pgfpathlineto{\pgfqpoint{0.908223in}{2.672228in}}%
\pgfpathlineto{\pgfqpoint{0.909614in}{2.767931in}}%
\pgfpathlineto{\pgfqpoint{0.911005in}{2.708090in}}%
\pgfpathlineto{\pgfqpoint{0.912396in}{2.717033in}}%
\pgfpathlineto{\pgfqpoint{0.913786in}{2.740430in}}%
\pgfpathlineto{\pgfqpoint{0.915177in}{2.851064in}}%
\pgfpathlineto{\pgfqpoint{0.916568in}{2.846476in}}%
\pgfpathlineto{\pgfqpoint{0.917959in}{2.813983in}}%
\pgfpathlineto{\pgfqpoint{0.919350in}{2.820888in}}%
\pgfpathlineto{\pgfqpoint{0.922132in}{2.784631in}}%
\pgfpathlineto{\pgfqpoint{0.923523in}{2.818867in}}%
\pgfpathlineto{\pgfqpoint{0.926305in}{2.692452in}}%
\pgfpathlineto{\pgfqpoint{0.927696in}{2.786497in}}%
\pgfpathlineto{\pgfqpoint{0.930478in}{2.744743in}}%
\pgfpathlineto{\pgfqpoint{0.933260in}{2.630227in}}%
\pgfpathlineto{\pgfqpoint{0.936042in}{2.660035in}}%
\pgfpathlineto{\pgfqpoint{0.938824in}{2.570159in}}%
\pgfpathlineto{\pgfqpoint{0.940215in}{2.594196in}}%
\pgfpathlineto{\pgfqpoint{0.941606in}{2.574796in}}%
\pgfpathlineto{\pgfqpoint{0.942997in}{2.582776in}}%
\pgfpathlineto{\pgfqpoint{0.944388in}{2.582960in}}%
\pgfpathlineto{\pgfqpoint{0.945779in}{2.558169in}}%
\pgfpathlineto{\pgfqpoint{0.947170in}{2.555752in}}%
\pgfpathlineto{\pgfqpoint{0.948560in}{2.517319in}}%
\pgfpathlineto{\pgfqpoint{0.949951in}{2.522313in}}%
\pgfpathlineto{\pgfqpoint{0.951342in}{2.495871in}}%
\pgfpathlineto{\pgfqpoint{0.954124in}{2.574368in}}%
\pgfpathlineto{\pgfqpoint{0.955515in}{2.560776in}}%
\pgfpathlineto{\pgfqpoint{0.958297in}{2.500869in}}%
\pgfpathlineto{\pgfqpoint{0.959688in}{2.568222in}}%
\pgfpathlineto{\pgfqpoint{0.961079in}{2.560749in}}%
\pgfpathlineto{\pgfqpoint{0.965252in}{2.644842in}}%
\pgfpathlineto{\pgfqpoint{0.966643in}{2.604857in}}%
\pgfpathlineto{\pgfqpoint{0.969425in}{2.648869in}}%
\pgfpathlineto{\pgfqpoint{0.970816in}{2.627969in}}%
\pgfpathlineto{\pgfqpoint{0.976380in}{2.818656in}}%
\pgfpathlineto{\pgfqpoint{0.977771in}{2.825920in}}%
\pgfpathlineto{\pgfqpoint{0.979162in}{2.848930in}}%
\pgfpathlineto{\pgfqpoint{0.980553in}{2.940452in}}%
\pgfpathlineto{\pgfqpoint{0.981944in}{2.949111in}}%
\pgfpathlineto{\pgfqpoint{0.983334in}{2.911299in}}%
\pgfpathlineto{\pgfqpoint{0.984725in}{2.928835in}}%
\pgfpathlineto{\pgfqpoint{0.986116in}{3.043829in}}%
\pgfpathlineto{\pgfqpoint{0.987507in}{3.066474in}}%
\pgfpathlineto{\pgfqpoint{0.988898in}{3.140910in}}%
\pgfpathlineto{\pgfqpoint{0.990289in}{3.102915in}}%
\pgfpathlineto{\pgfqpoint{0.994462in}{3.142907in}}%
\pgfpathlineto{\pgfqpoint{0.995853in}{3.108715in}}%
\pgfpathlineto{\pgfqpoint{0.998635in}{3.003748in}}%
\pgfpathlineto{\pgfqpoint{1.000026in}{3.033466in}}%
\pgfpathlineto{\pgfqpoint{1.001417in}{3.038711in}}%
\pgfpathlineto{\pgfqpoint{1.002808in}{3.024787in}}%
\pgfpathlineto{\pgfqpoint{1.004199in}{3.169487in}}%
\pgfpathlineto{\pgfqpoint{1.005590in}{3.179014in}}%
\pgfpathlineto{\pgfqpoint{1.006981in}{3.214198in}}%
\pgfpathlineto{\pgfqpoint{5.892034in}{3.214198in}}%
\pgfpathlineto{\pgfqpoint{5.892034in}{3.214198in}}%
\pgfusepath{stroke}%
\end{pgfscope}%
\begin{pgfscope}%
\pgfpathrectangle{\pgfqpoint{0.050000in}{0.050000in}}{\pgfqpoint{6.120226in}{3.883333in}}%
\pgfusepath{clip}%
\pgfsetbuttcap%
\pgfsetroundjoin%
\pgfsetlinewidth{1.003750pt}%
\definecolor{currentstroke}{rgb}{0.000000,0.000000,0.000000}%
\pgfsetstrokecolor{currentstroke}%
\pgfsetdash{{1.000000pt}{1.650000pt}}{0.000000pt}%
\pgfpathmoveto{\pgfqpoint{1.006981in}{1.775926in}}%
\pgfpathlineto{\pgfqpoint{1.006981in}{3.214198in}}%
\pgfusepath{stroke}%
\end{pgfscope}%
\begin{pgfscope}%
\pgfsetbuttcap%
\pgfsetmiterjoin%
\definecolor{currentfill}{rgb}{0.000000,0.000000,0.000000}%
\pgfsetfillcolor{currentfill}%
\pgfsetlinewidth{1.003750pt}%
\definecolor{currentstroke}{rgb}{0.000000,0.000000,0.000000}%
\pgfsetstrokecolor{currentstroke}%
\pgfsetdash{}{0pt}%
\pgfsys@defobject{currentmarker}{\pgfqpoint{-0.027778in}{-0.027778in}}{\pgfqpoint{0.027778in}{0.027778in}}{%
\pgfpathmoveto{\pgfqpoint{0.027778in}{-0.000000in}}%
\pgfpathlineto{\pgfqpoint{-0.027778in}{0.027778in}}%
\pgfpathlineto{\pgfqpoint{-0.027778in}{-0.027778in}}%
\pgfpathlineto{\pgfqpoint{0.027778in}{-0.000000in}}%
\pgfpathclose%
\pgfusepath{stroke,fill}%
}%
\begin{pgfscope}%
\pgfsys@transformshift{6.170226in}{1.775926in}%
\pgfsys@useobject{currentmarker}{}%
\end{pgfscope}%
\end{pgfscope}%
\begin{pgfscope}%
\pgfsetbuttcap%
\pgfsetmiterjoin%
\definecolor{currentfill}{rgb}{0.000000,0.000000,0.000000}%
\pgfsetfillcolor{currentfill}%
\pgfsetlinewidth{1.003750pt}%
\definecolor{currentstroke}{rgb}{0.000000,0.000000,0.000000}%
\pgfsetstrokecolor{currentstroke}%
\pgfsetdash{}{0pt}%
\pgfsys@defobject{currentmarker}{\pgfqpoint{-0.027778in}{-0.027778in}}{\pgfqpoint{0.027778in}{0.027778in}}{%
\pgfpathmoveto{\pgfqpoint{0.000000in}{0.027778in}}%
\pgfpathlineto{\pgfqpoint{-0.027778in}{-0.027778in}}%
\pgfpathlineto{\pgfqpoint{0.027778in}{-0.027778in}}%
\pgfpathlineto{\pgfqpoint{0.000000in}{0.027778in}}%
\pgfpathclose%
\pgfusepath{stroke,fill}%
}%
\begin{pgfscope}%
\pgfsys@transformshift{0.328192in}{3.933333in}%
\pgfsys@useobject{currentmarker}{}%
\end{pgfscope}%
\end{pgfscope}%
\begin{pgfscope}%
\pgfsetrectcap%
\pgfsetmiterjoin%
\pgfsetlinewidth{0.501875pt}%
\definecolor{currentstroke}{rgb}{0.000000,0.000000,0.000000}%
\pgfsetstrokecolor{currentstroke}%
\pgfsetdash{}{0pt}%
\pgfpathmoveto{\pgfqpoint{0.328192in}{0.050000in}}%
\pgfpathlineto{\pgfqpoint{0.328192in}{3.933333in}}%
\pgfusepath{stroke}%
\end{pgfscope}%
\begin{pgfscope}%
\pgfsetrectcap%
\pgfsetmiterjoin%
\pgfsetlinewidth{0.000000pt}%
\definecolor{currentstroke}{rgb}{0.000000,0.000000,0.000000}%
\pgfsetstrokecolor{currentstroke}%
\pgfsetstrokeopacity{0.000000}%
\pgfsetdash{}{0pt}%
\pgfpathmoveto{\pgfqpoint{6.170226in}{0.050000in}}%
\pgfpathlineto{\pgfqpoint{6.170226in}{3.933333in}}%
\pgfusepath{}%
\end{pgfscope}%
\begin{pgfscope}%
\pgfsetrectcap%
\pgfsetmiterjoin%
\pgfsetlinewidth{0.501875pt}%
\definecolor{currentstroke}{rgb}{0.000000,0.000000,0.000000}%
\pgfsetstrokecolor{currentstroke}%
\pgfsetdash{}{0pt}%
\pgfpathmoveto{\pgfqpoint{0.050000in}{1.775926in}}%
\pgfpathlineto{\pgfqpoint{6.170226in}{1.775926in}}%
\pgfusepath{stroke}%
\end{pgfscope}%
\begin{pgfscope}%
\pgfsetrectcap%
\pgfsetmiterjoin%
\pgfsetlinewidth{0.000000pt}%
\definecolor{currentstroke}{rgb}{0.000000,0.000000,0.000000}%
\pgfsetstrokecolor{currentstroke}%
\pgfsetstrokeopacity{0.000000}%
\pgfsetdash{}{0pt}%
\pgfpathmoveto{\pgfqpoint{0.050000in}{3.933333in}}%
\pgfpathlineto{\pgfqpoint{6.170226in}{3.933333in}}%
\pgfusepath{}%
\end{pgfscope}%
\end{pgfpicture}%
\makeatother%
\endgroup%

	\caption[Stopped stochastic process]{A Stopped Stochastic Process}\label{fig:StoppedStochasticProcess}
\end{figure}

Suppose that a martingale is stopped (frozen) if it will go up in the near future, but is allowed to continue if it will go down. Then, stopping introduces a downward bias by removing the upward possibility.

\begin{prop}{Stochastic Integral for Predictable Integrands}{StochasticIntegralForPredictableIntegrands}
	Let \(N\) be defined as above. Then, for any positive \(\cF\)-predictable process \((F_t)\),
	\begin{equation}\label{eq:StochasticIntegralExpectation}
		\E \int_{\R_+} F_t \di N_t = \E \int_{\R_+} F_t \di \nu_t.
	\end{equation}
\end{prop}

\begin{thm}{Integral Transform of a Continuous Martingale}{IntegralTransformOfAContinuousMartingale}
	Let \(\overline{N}\) be an \(\cF\)-martingale, as defined above. For any bounded \(\cF\)-predictable process \(F\), the integral transform
	\begin{equation}\label{eq:ContinuousMartingaleTransform}
		Z_t \coloneqq \int_0^t F_s \di \overline{N}_s, \quad t \ge 0,
	\end{equation}
	is a martingale.
\end{thm}

\subsubsection{Doob's Stopping Theorem}
\begin{thm}{Doob's Stopping Theorem}{DoobsStoppingTheorem}
	A stochastic process \(M \coloneqq (M_n, n \in \N)\) is a martingale if and only if for every pair of \emph{bounded} stopping times \(S\) and \(T\) with \(S \le T\) the random variables \(M_S\) and \(M_T\) are integrable, and
	\begin{equation}\label{eq:DoobsStoppingEquality}
		\E_S(M_T - M_S) = 0.
	\end{equation}
\end{thm}

\begin{lem}{Maximum Inequality for Positive Submartingales}{MaximumInequalityForPositiveSubmartingales}
	Let \((Y_k, k \in \N)\) be a positive submartingale and define \(Y_n^* \coloneqq \max_{k \le n} Y_k\). Then, for any \(b > 0\),
	\begin{equation}\label{eq:MaxInequalityPositiveSub}
		b \Pm(Y_n^* \ge b) \le \E\bigl[Y_n \I_{\{Y_n^* \ge b\}}\bigr] \le \E Y_n.
	\end{equation}
\end{lem}

\begin{thm}{Doob--Kolmogorov Inequality}{DoobKolmogorovInequality}% chktex 8
	Let \(M \coloneqq (M_k)\) be a martingale in \(L^p\) for some \(p\) with \(1 \le p < \infty\). Then, for \(b > 0\),
	\begin{equation}\label{eq:DoobKolmogorovInequality}
		\Pm\bigl(\max_{k \le n} \lvert M_k \rvert > b\bigr) \le \frac{\E \lvert M_n \rvert^p}{b^p}.
	\end{equation}
\end{thm}

\subsection{(Sub) Martingale Convergence}
The second major use of martingales is in the proof of convergence for stochastic processes. A first step in the analysis is to characterise the number of times a stochastic process ``upcrosses'' an interval.

\subsubsection{Upcrossing}

Let \(X \idef (X_k, k \in \N)\) be an adapted process with respect to some filtration and consider the integral transformation \(Z_n \idef \int_0^n F_k \di X_k\), \(n \in \N\), where \((F_k)\) is a predictable process. Continuing the financial analogy, think of \(X_k\) as the share price at time \(k\) and \(F_k\) as the number of shares owned at period \(k\). It is completely determined by the share price history before period \(k\). That is, the share buying strategy is determined by the information available before the \(k\)th period.

For example, the strategy could have the following rules: (1) we are allowed at most one share, (2) when the share price drops to level \(a\) or lower, we buy one share (if our portfolio is empty), and (3) when the share price reaches \(b\) or higher, for some fixed \(b > a\), we sell the share, if we own one.
The total profit at period \(n\) is \(Z_n - Z_0\).
If the share price at \(n\) (if any) is higher than what it was bought for, then the total profit \(Z_n - Z_0\) is at least \((b - a) U_n\), where \(U_n\) is the number of times the process \(X\) has \textit{upcrossed} the interval \((a, b)\) during the periods \(0, 1, \ldots, n\).

As a special case, suppose that \(X \ge 0\) is a \textit{positive submartingale} and that \(a = 0\). Thus, if at any time \(k\) when the price drops to \(0\), we buy a share (if we have no share), and we sell a share (if we own one) when the price hits \(b\) or any price higher. Again, we are allowed at most one share. We thus have
\begin{equation}\label{eq:SpecialCaseUpcrossingProfit}
	Z_n - Z_0 \ge b \, U_n,
\end{equation}
where \(U_n\) is the number of upcrossings of \((0, b)\); i.e., the total number of shares sold up to period \(n\).

Now observe that for each \(k = 0, 1, 2, \ldots\), we have:
\begin{equation}\label{eq:PredictableIncrementBound}
	\E_k[Z_{k+1} - Z_k]
	= \E_k[(X_{k+1} - X_k) F_{k+1}]
	= \underbrace{F_{k+1}}_{\le 1} \, \underbrace{\E_k[X_{k+1} - X_k]}_{\ge 0}
	\le \E_k[X_{k+1} - X_k].
\end{equation}

Summing the above terms from \(k = 0\) to \(n - 1\) and taking expectations gives
\begin{equation}\label{eq:ExpectedIncrementBound}
	\E[Z_n - Z_0] \le \E[X_n - X_0].
\end{equation}

Combining this with \Cref{eq:SpecialCaseUpcrossingProfit} gives:
\begin{equation}\label{eq:ExpectedUpcrossingsSpecialCase}
	b \, \E U_n \le \E[X_n - X_0].
\end{equation}


For a general interval \((a, b)\) and a submartingale \(X\), the above result is generalized as follows. Recall that for a random variable \(V\), \(V^{+} \idef V \lor 0 = \max(V, 0)\).


\begin{thm}{Expected Upcrossings of a Submartingale}{ExpectedUpcrossingsOfASubmartingale}
	Suppose that \(X\) is a submartingale. Then, the expected number of upcrossings \(U_n(a, b)\) of \((a, b)\) satisfies:
	\begin{equation}\label{eq:UpcrossingsInequality}
		(b - a) \E U_n(a, b) \le \E\bigl[{(X_n - a)}^{+} - {(X_0 - a)}^{+}\bigr].
	\end{equation}
\end{thm}

\begin{proof}
	The number of upcrossings of \((a, b)\) for \(X\) is the same as the number of upcrossings of \((0, b - a)\) for the process \({(X - a)}^{+}\). The latter is a positive submartingale, for which, by \Cref{eq:ExpectedUpcrossingsSpecialCase}, we have \Cref{eq:UpcrossingsInequality}.
\end{proof}

The upcrossing theorem will be useful for proving convergence results for (sub) martingales.
Recall that a submartingale has a tendency to increase.
If this increase is constrained in some way, it is reasonable that the process would converge.
The following theorem restrains the growth of the submartingale by means of a finiteness condition on the supremum of the \(\E X_n^{+}\):

\begin{thm}{(Sub) Martingale Convergence}{SubMartingaleConvergence}
	Let \(X \coloneqq (X_n, n \in \N)\) be a submartingale. Suppose that \(\sup_n \E X_n^+ < \infty\). Then, \(X\) converges almost surely to an integrable random variable \(X_\infty\).
\end{thm}

\begin{thm}{(Sub) Martingale Convergence and Uniform Integrability}{SubMartingaleConvergenceAndUniformIntegrability}
	Let \(X \coloneqq (X_n, n \in \N)\) be a submartingale. Then, \(X\) converges almost surely and in \(L^1\) if and only if it is uniformly integrable. Moreover, if \(X\) converges, setting \(X_\infty \coloneqq \lim X_n\) extends \(X\) to a submartingale \(\overline{X} \coloneqq (X_n, n \in \overline{\N})\).
\end{thm}

\begin{thm}{Uniformly Integrable Martingale}{UniformlyIntegrableMartingale}
	A process \(M \coloneqq (M_n, n \in \N)\) is a UI martingale if and only if there exists an integrable random variable \(Z\) such that
	\begin{equation}\label{eq:UIMartingaleDef}
		M_n = \E_n Z, \quad n \in \N.
	\end{equation}
	Moreover, then \(M\) converges almost surely and in \(L^1\) to an integrable random variable
	\begin{equation}\label{eq:MartingaleLimit}
		M_\infty \coloneqq \E_\infty Z,
	\end{equation}
	and \((M_n, n \in \overline{\N})\) is again a UI martingale. If \(Z \in \cF_\infty\), then \(M_\infty = Z\).
\end{thm}

\begin{thm}{Convergence for Martingales in Reversed Time}{ConvergenceForMartingalesInReversedTime}
	For \(\bb{T} = \{\dots, -2, -1, 0\}\), let \(M \coloneqq (M_n, n \in \bb{T})\) be a martingale with respect to the filtration \((\cF_n, n \in \bb{T})\). Then, \(M\) is UI and, moreover, as \(n \downarrow -\infty\) it converges almost surely and in \(L^1\) to the integrable random variable \(M_{-\infty} \coloneqq \E_{-\infty} M_0\), where \(\cF_{-\infty} \coloneqq \cap_{n \in \bb{T}} \cF_n\).
\end{thm}

\begin{thm}{Stopping for Uniformly Integrable Martingales}{StoppingForUniformlyIntegrableMartingales}
	If \(M \coloneqq (M_n, n \in \overline{\N})\) is a uniformly integrable martingale, then for \emph{every} pair of stopping times \(S\) and \(T\) with \(S \le T\) it holds that
	\begin{equation*}
		\E_S M_T = M_S
	\end{equation*}
	In particular,
	\begin{equation*}
		\E M_T = \E M_0.
	\end{equation*}
\end{thm}

\subsection{Applications}

\subsubsection{Kolmogorov's 0--1 Law}
\begin{thm}{Kolmogorov's 0--1 Law}{KolmogorovsZeroOneLaw}
	If \(H \in \mathcal{T}\), then \(\Pm(H)\) is either \(0\) or \(1\). Consequently, a numerical random variable \(Z\) that is \(\mathcal{T}\)-measurable must be almost surely constant.
\end{thm}

\subsubsection{Radon--Nikodym Theorem}% chktex 8
\begin{defn}{Separable \(\sigma\)-Algebra}{SeparableSigmaAlgebra}
	A \(\sigma\)-algebra on \(\Omega\) is said to be \emph{separable} if it is generated by a sequence \((H_n)\) of subsets of \(\Omega\).
\end{defn}

\begin{thm}{Radon--Nikodym}{RadonNikodym}% chktex 8
	Let \((\Omega, \cH, \Pm)\) be a probability space, where \(\cH\) is separable, and let \(Q\) be a finite measure on \((\Omega, \cH)\) with \(Q \ll \Pm\). Then, there exists a positive random variable \(Z \in \cH\), written \(Z = \di Q / \di \Pm\), such that \(Q = Z \Pm\); that is,
	\begin{equation}\label{eq:RadonNikodymFormula}
		Q(H) = \int_H Z \di \Pm = \E \I_H Z, \quad H \in \cH.
	\end{equation}
\end{thm}

\begin{defn}{Right-continuous Filtration}{RightContinuousFiltration}
	A filtration \(\cF \coloneqq (\cF_t)\) is said to be \emph{right-continuous} if for all \(t\),
	\begin{equation*}
		\bigcap_{\varepsilon > 0} \cF_{t+\varepsilon} = \cF_t.
	\end{equation*}
\end{defn}

\begin{defn}{Local Martingale}{LocalMartingale}
	A process \((Z_t, t \ge 0)\) is called a \emph{local martingale} if there exists a sequence of stopping times \((T_n, n \in \N)\), called a \emph{localizing sequence}, such that \(T_n \to \infty\) almost surely and each stopped process \((Z_{t \land T_n}, t \ge 0)\) is a martingale.
\end{defn}

\begin{thm}{Properties of Local Martingales}{PropertiesOfLocalMartingales}
	Let \(Z \coloneqq (Z_t, t \ge 0)\) be a local martingale with localizing sequence \((T_n, n \in \N)\). Then, the following hold:
	\begin{enumerate}
		\item If \(Z \ge 0\) and \(\E Z_0 < \infty\), then \(Z\) is a supermartingale.
		\item If for each \(t \ge 0\) the sequence \((Z_{t \land T_n}, n \in \N)\) is uniformly integrable, then \(Z\) is a martingale.
	\end{enumerate}
\end{thm}

\begin{defn}{Doob Martingale}{DoobMartingale}
	Let \(\zeta\) be a stopping time. A process \(M \coloneqq (M_t, t \in \bb{T})\) is called a \emph{Doob martingale} on \([0, \zeta]\) if for all stopping times \(0 \le S \le T \le \zeta\) it holds that
	\begin{equation}\label{eq:DoobMartingaleCondition}
		\E_S M_T = M_S,
	\end{equation}
	where \(M_S\) and \(M_T\) are integrable.
\end{defn}

\begin{thm}{Characterization of a Doob Martingale}{CharacterizationOfADoobMartingale}
	Let \(\zeta\) be a stopping time. A process \(M \coloneqq (M_t, t \in \bb{T})\) is a Doob martingale on \([0, \zeta]\) if and only if for every stopping time \(T \le \zeta\),
	\begin{equation}\label{eq:DoobMartingaleExpectation}
		\E M_T = \E M_0.
	\end{equation}
\end{thm}

\begin{thm}{Doob's Stopping Theorem for Continuous Martingales}{DoobsStoppingTheoremForContinuousMartingales}
	A martingale \(M \coloneqq (M_t, t \in \R_+)\) is a Doob martingale on \([0, b]\) for every \(b \in \R_+\) and satisfies
	\begin{equation}\label{eq:DoobStoppingContinuous}
		\E_S M_T = M_S
	\end{equation}
	for every pair \((S, T)\) of bounded stopping times with \(S \le T\).
\end{thm}

\begin{prop}{Doob's Maximal Inequality}{DoobsMaximalInequality}
	Let \(X \coloneqq (X_t, t \ge 0)\) be a submartingale that is positive and continuous. Then, for \(p \ge 1\) and \(b > 0\),
	\begin{equation}\label{eq:DoobsMaximalInequality}
		\Pm(\max_{0 \le s \le t} X_s \ge b) \le \frac{\E X_t^p}{b^p}.
	\end{equation}
\end{prop}

\begin{prop}{Doob's Norm Inequality}{DoobsNormInequality}
	Let \(M \coloneqq (M_t, t \ge 0)\) be a continuous martingale in \(L^p\) for some \(p > 1\) and let \(Z_t \coloneqq \max_{s \le t} \lvert M_s \rvert\). Then,
	\begin{equation}\label{eq:DoobsNormInequality}
		\E Z_t^p \le q^p \E \lvert M_t \rvert^p, \text{ with } q \coloneqq p/(p-1).
	\end{equation}
\end{prop}

\begin{lem}{Continuous Martingales with Finite Variation}{ContinuousMartingalesWithFiniteVariation}
	Let \(X \coloneqq (X_t, t \ge 0)\) be a martingale with continuous paths and total variation \(V_t < \infty\) on each interval \([0, t]\). Then, \(X\) is almost surely constant.
\end{lem}

\begin{prop}{Doob Martingale on \(\overline{\R}_+\)}{DoobMartingaleOnRBar}
	A process \(M\) is a Doob martingale on \(\overline{\R}_+\) if and only if it is a martingale on \(\overline{\R}_+\). If so, then
	\begin{equation}\label{eq:DoobMartingaleRBarProp}
		\E_S M_T = M_S
	\end{equation}
	for arbitrary stopping times with \(S \le T\).
\end{prop}

\begin{thm}{Extension of UI Martingales}{ExtensionOfUIMartingales}
	A martingale \(M\) on \(\R_+\) can be extended to a Doob martingale \(\overline{M}\) on \(\overline{\R}_+\) if and only if it is uniformly integrable; that is, if and only if
	\begin{equation}\label{eq:UIMartingaleExtension}
		M_t = \E_t Z, \quad t \in \R_+
	\end{equation}
	for some integrable random variable \(Z\). Moreover, then, it converges almost surely and in \(L^1\) to an integrable random variable \(M_\infty\) and can be extended to a UI martingale on \(\overline{\R}_+\).
\end{thm}

\begin{prop}{Sufficient Condition for Being Doob}{SufficientConditionForBeingDoob}
	Let \(M\) be a martingale on \(\R_+\) and \(\zeta\) a stopping time. If, almost surely,
	\begin{equation}\label{eq:DoobDominationCondition}
		\sup_{t \in [0,\zeta] \cap \R_{+}} \lvert M_t \rvert \le Z,
	\end{equation}
	where \(\E Z < \infty\), i.e., \(M\) is dominated by an integrable random variable \(Z\). Then \(M_\zeta = \lim_{t \to \infty} M_{\zeta \land t}\) exists and is integrable. 	Moreover, \(M\) is a Doob martingale on \([0, \zeta]\).
\end{prop}% chktex 17
